\subsection{高分母 support の現在地と，適用仮定}
\label{ns19:current-support}

対象は，全平面五次の標準 $K_5$ branch graph と $K_4$ cones を保つ作用適合 atlas である．
$H=\{ij:q_{ij}>1\}$ とする．木・単純閉路成分，二価二頂点を持つ三角形，
葉辺付き四角形，$K_{2,3}$，diamond，二葉付き三角形，四使用頂点上の $K_4$ に加え，
$G_{5,5},G_{5,10},G_{5,11},G_{5,12},G_{5,13}$ と葉付き $G_{5,7},G_{5,9}$ を排除した．これらの障害は分母の数値上限を用いない．
現在の必要候補は\emph{三型・二十六標識}であり，整数接着や atlas の存在を主張しない．

非空 $H$ は五頂点全てを使う三点連結 graph で，最小次数三以上，少なくとも八辺を持つ．
二価 core を持つ最後の候補 $G_{5,15}$ も，通常型と負例外の全尺度排除により解決した．
一つの core に高分母 ray が高々二本なら全分母一である．
従って非空 $H$ の単位枝は高々二辺の matching となる．
これは現在の graph に対する結論であり，別の core や branch singularity を導入する場合には，
変更後の接着条件を定義してから障害の適用範囲を調べる．

\begin{center}
\begin{tabular}{lcc}
障害 & 対向 $D$ の相殺 & node 非負性\\\hline
木成分・三角形成分 & 用いる & 不要\\
四角形・五角形成分 & 不要 & 用いる\\
二価二頂点を持つ三角形 & 不要 & 用いる\\
四角形と一本の葉辺 & 用いる & 用いる\\
$K_{2,3}$・負の枝の三角形条件 & 不要 & 用いる\\
diamond・二葉付き三角形の全排除 & 用いる & 用いる\\
全 support が $K_4$ & 低 hub 辺で用いる & 用いる\\
$G_{5,5}$ & 二本の低辺で用いる & 用いる\\
$G_{5,10}$ & 用いる & 用いる\\
$G_{5,12}=K_4+\mathrm{leaf}$ & 葉の高辺で用いる & 用いる\\
$G_{5,7}$ & 葉の高辺で用いる & 用いる\\
$G_{5,9}$ & 葉に接する四辺で用いる & 用いる\\
$G_{5,11}$・book $G_{5,13}$ & 低辺で用いる & 用いる
\end{tabular}
\end{center}
node 全数五および各 node の上限五は row sums と非負性から導く．
三角形の任意の追加高辺を許したわけではなく，二頂点が二価という仮定を維持する．
以下には，今後も使える定量的な伝播式を残す．

\subsection{三高分母の完全マッチング定理：今後は仮定でなく入力として使う}
\label{nm18:matching}

整数行和・pair products を満たし，$q_0=1,q_1,q_2,q_3>1$ である局所模型を考える．
正の entries の無向 graph は，完全マッチングまたは低 ray を含む三角形と孤立頂点の和に限られる．
さらに，正の basin weights を持つことと，完全マッチング型であることは同値である．
分母上限，clutching，node，対向 $D$ 条件はいずれも不要である．
証明は算術稿の \texttt{mh18:matching} に置く．

再利用時には，低 row の三つの entries がちょうど二負一正であることを最初に課す．
正の低 entry の相手を $c$，負の二つの相手を $a,b$ とすると，
正符号 graph は $0c\sqcup ab$ であり，既存の canonical Schur 表示が全ての正の三高模型を覆う．
$\beta\le-1$ と，境界 $\beta=-1$ の完全な約数表示をそのまま使える．
ただし，全ての高 $D$ が非負となる可能性はまだ解決していない．
この係数問題の残りは $b=-\beta\ge2$，$d=-D_0\ge1$，
$\alpha\ge k\ge b+1$ の非特異領域である．
特異 Schur 系は全高 $D$ 非負の候補から除かれたが，
一般の正整数局所模型の分類にはその実在する無限族を保持する．

\paragraph{他の問題にも使える二つの証明操作．}
二つの道型では，同じ残余二次式
\[
 (kC-b)L^2+\{kd-aC+B(kC-b)\}L+d(kB-a)=0
\]
に対し，全係数の符号がそろう．正の根 $L$ がないことが，無上界の排除を与える．
三角形型では，正の三閉路の不等式を
\[
 (tS-H)K<1,\qquad (sT-K)H<1
\]
と書くと，整数性から $H=tS+\alpha,K=sT+\beta$，$\alpha,\beta\ge0$ となる．
残る pair equations を整理した
\[
 bS=AT+\alpha\beta,\qquad
 (A-1)U=CS+k(\alpha+1),
\]
\[
 (b+1)U=(C+1)T+(k+1)(\beta+1)
\]
は $A>1,b>0,U>T$ を強制する．$\delta=U-T\ge1$ に対して
\[
 A\{s(1-\delta)-k\beta\}=bsU+bk(s+t)+Ct\beta
\]
の左辺は非正，右辺は正となる．
ここでは $A=s\alpha-t$，$b=s-t\beta$，$C=s+(s+t+1)k$ であり，
この排除に用いる変数は canonical matching 型の変数と混同しない．
実数での $U>T$ を整数の $\delta\ge1$ へ強める一歩が不可欠である．

\paragraph{正値性を省く誤りを避ける検査例．}
分母 $(1,5,19,31)$ と
\[
 P=\begin{pmatrix}
 0&1&1&-3\\4&0&2&-11\\18&47&0&-84\\-10&-14&-7&0
 \end{pmatrix}
\]
は整数行和・pair products を全て満たし，正符号 graph は三角形である．
しかし閉路 $0\to2\to1\to0$ の gain は $188/95>1$ であり，正 weights はない．
行和・pair products だけから完全マッチングを結論してはいけない．

\nolinkurl{verify_three_high_matching_v18.py} は十一の多項式恒等式を係数ごとに確認し，
道型の係数の符号を正変数の平行移動で検査する．
併せて固定低データの小さい入力を既存の完全復元器へ通し，実装の整合性を点検する．
この入力格子の検査を，一般定理の証明として引用しない．


\paragraph{別の接着問題にも再利用する葉の伝播則．}
葉 $i$，二価 core $j$，その先の core $k$ を考える．両方の高辺を実際に接続すると，
ある $h\geq1$ により
\[
 a=q_{ij}=h(h+1)+1,\qquad b=q_{jk}=(h+1)(h+2)+1,
\]
\[
 D^j_i=a,\qquad D^j_k=-b,\qquad
 p^k_{ji}=-(h+1),\qquad p^k_{ij}=-(h+2)
\]
に強制される．特に $b>a$，$\gcd(a,b)=1$ である．
最初の分母対は $(3,7),(7,13),(13,21),(21,31)$ となる．
この結果は $k$ の高分母次数を仮定せず，node 条件も用いないため，
cycle を含む成分に付いた長さ二の pendant path にも適用できる．
現行の四候補では葉自体が存在しないが，この式は node 条件を緩めた接着問題にも使える．
ただしこの必要条件だけで接着可能性を結論しない．

第三 core の条件を省くと，二つの追加族
\[
 (a,b,D^j_i,D^j_k)=(2t+1,6t+1,3t,1-8t),
\]
\[
 (a,b,D^j_i,D^j_k)=(6t+1,20t-1,7t,3-18t)
\]
が残る．後者は $D^j_k>-b$ なので，一般の二高分母の局所不等式に移し替えてはならない．
ところが次の core から元の葉へ向かう entry は $b-1$ の signed divisor で，
前者では $3t+2\pmod{6t+1}$，後者では $10t+2\pmod{20t-1}$ を満たす必要がある．
前者は $t=1$ だけ，後者は全て不可能となる．
一本の辺での整合と，次の辺までの整合の差を示す再利用可能な例である．

\paragraph{木の中心で使う係数和の上界．}
三高分母 $(a,b,c,1)$ の分母一の vertex を $0$ とし，
$r_i=p_{0i},x_i=p_{i0}$ と置く．整数性と局所式だけで
\[
 \sum r_i=-1,\qquad r_ix_i=q_i-1>0,\qquad
 \sum_{i=a,b,c}D_i=1+\sum x_i
\]
である．ある $x_j\leq-1$ が存在し，他の $x_i\leq q_i-1$ なので
\[
 D_a+D_b+D_c\leq a+b+c-\min\{a,b,c\}-2.
\]
一高分母の葉は反対側に $D\geq q$ を要求し，上記の二価 core を経由しても同じ下界が届く．
従って三価の star とその一辺を細分した木は，中心でこの上界に反する．
四価 star は全 $D$ の和零で排除される．
正局所模型では完全 matching 定理により，low row とその逆向き列は二負一正である．
正の entry が高 ray $c$ に対応すれば $x_a,x_b\le-1$，$x_c\le q_c-1$ なので，さらに
\[
 D_a+D_b+D_c=1+x_a+x_b+x_c\le q_c-2\le\max\{a,b,c\}-2
\]
を得る．これは新しい matching 性から直ちに得る上界であり，
一般三高分母の「全高 $D>0$ 不可」を仮定するものではない．

\paragraph{四角形と五角形で再利用する node の使い方．}
二高分母の負の高低 entry は $-(q-1)$ にはならず，絶対値は $(q-1)/2$ 以下である．
従って，共通の低 puncture で両端が負となることは整数 clutching と矛盾する．
四角形では正の pairing が全周で同じ向きになる．mixed node の二型を伝播させると，
一方の型は第五の全低分母 core により canonical $S=T=2$ を強制する．
この局所型は $(r,s,z,a,b)=(4,4,5,9,11)$ だけなので，四つの core の共有分母が一致しない．
他方の型は逆向き約数 $(4,3,2,1)$ と真約数性が矛盾する．
この改訂証明は $D$ を使わない．

五角形では，正の pairing の vertex は高々一つである．残る四つの連続した負の pairing から，
三本の高辺の共通低 puncture は全て正となり，それぞれの両端の値は $\{1,q-1\}$ となる．
canonical な $A$ だけが正の最大約数 $q-1$ を持ち得るため，これらの端が伝播する．
その結果，ある低辺の node 数が $n=2-q<0$ となる．
この三本の例外辺の伝播は，整数 clutching と node 条件を同時に課す手順として再利用できる．

\subsection{二高分母の最小 entry と，誤った強化への反例}

canonical な非零 low--low entry $d=S-T-1$ は，$\alpha,\beta<0$ のときちょうど $-1$，
$\alpha,\beta>0$ なら $d\leq-3$ となり，$d=-2$ はない．
$d=-1$ では大きい高分母の $D$ が負，小さい方が正である．
ただし，大域標識で固定した向きの low--low entry は零の場合もある．

この最小 entry 型にも，正の高係数和を持つ無限族がある．$H\geq2$ として
\[
 r=2H,\quad S=T=H,\quad s=2H,\quad z=(2H+1)(H-1)
\]
と置くと，
\[
 a=2H^2+1,\qquad b=H(H-1)(2H+1)+1,
\]
\[
 D_A=2H^3-2H^2-2H-1>0,\qquad
 D_B=-2H^3+3H^2+H+1<0
\]
であり，$D_A+D_B=H(H-1)>0$ となる．
全 row sums と pair products は恒等式であり，明示 weights も存在する．
従って，高係数和の一様な負性を再び仮定しない．

\subsection{固定分母の相手と，小分母の係数を保持した伝播}

奇分母 $n>1$ を含む二高分母模型の相手集合 $\mathcal P(n)$ は，$n\ne7,13$ なら有限である．
七と十三だけが共鳴し，それぞれ $6k+1$ と $12k+1$，$k\geq1$，という無限族を含む．
これら以外の相手は有限個である．canonical な $A/B$ の役割を分けることが不可欠である．
$A$ 側を固定すれば
\[
 z\{A(r+1)-(1+rs)T\}=(r+1)(s-x-1),
\]
$B$ 側なら
\[
 s\{(1+(r+1)z)S-rB\}=r(y-z-1)
\]
を用いる．右辺が零の枝を落とさずに扱う．
\nolinkurl{verify_fixed_denominator_partners_v12.py} の
\texttt{partners\_for\_fixed(n)} は，全非共鳴標準形と共鳴枝を返す．

相手の分母だけでなく固定 ray の $D$ を保存すると，二高分母 core 同士の接続について
次の必要条件を得る．共有分母三・五は不可能である．七では，非零係数なら $\pm7$ で，
隣接高分母は $\{3,13\}$ または $\{3,19\}$，零なら両側とも七である．
十三では，非零係数なら隣接高分母 $\{5,25\}$ と係数 $\pm15$，または
$\{7,21\}$ と係数 $\pm13$ に限られ，零なら両側の隣接分母は各々十三または十七である．
この零係数の場合を一般の接続から落としてはいけない．
固定七・十三の共鳴族の係数 $7k-9,13k-14$ は，もう一つの二高分母 core により相殺できない．
全四十七非共鳴模型と二共鳴のデータは v13 の小分母伝播検算に保存した．

% v13 research notes: reusable local classification, not an unresolved C5 claim.
\subsection{二つの負の reverse divisors が四以下の完全な局所 catalog}
\label{nr13:bounded-reverse}

二高分母標準形の負の高低 entries は $-s,-y$ であり，
それぞれの reverse numerators の絶対値は $S,r+1$ である．
これらがともに四以下という局所条件を課すと，
分母の上限を置かずに全ての整数模型を有限表と三つの族に分けられる．
この catalog は，今後別の support に現れる共通 $\lambda=0$ punctures を調べる際に再利用できる．
五角形成分自体は v13 で排除されており，その残存候補を意味しない．

\paragraph{仮定と完全性．}
正の整数二高分母標準形で $a,b$ は奇数とし，
\[
 2\leq S\leq4,\qquad 1\leq r\leq3
\]
を仮定する．$S\geq2$ は負の低分母 entry の真約数性から従う．
このとき，次の表の二十二模型と三つの無限族だけが可能である．
表は canonical data の分類であって，標識を忘れた幾何的同型類を数えるものではない．

\begin{center}
\small
\begin{tabular}{rrrrrrrrr}
$r$&$S$&$T$&$s$&$z$&$a$&$b$&$D_A$&$D_B$\\\hline
1&2&5&1&2&3&11&6&$-10$\\
1&2&5&2&6&5&31&18&$-27$\\
1&3&8&2&5&7&41&21&$-36$\\
1&3&9&2&2&7&19&7&$-15$\\
1&3&9&4&6&13&55&23&$-41$\\
1&3&9&6&26&19&235&127&$-211$\\
1&3&11&8&2&25&23&$-9$&9\\
1&3&12&4&1&13&13&$-3$&$-2$\\
1&4&9&1&4&5&37&21&$-36$\\
1&4&10&1&2&5&21&11&$-21$\\
1&4&12&1&1&5&13&6&$-15$\\
1&4&12&3&5&13&61&26&$-49$\\
1&4&12&4&13&17&157&78&$-138$\\
1&4&13&9&28&37&365&177&$-312$\\
1&4&13&10&62&41&807&428&$-747$\\
1&4&15&6&2&25&31&$-4$&$-1$\\
1&4&15&25&4&101&61&$-63$&102\\
1&4&16&3&1&13&17&0&$-9$\\
1&4&16&13&2&53&33&$-31$&45\\
1&4&18&7&1&29&19&$-15$&15\\
2&3&5&2&6&7&31&21&$-25$\\
2&4&7&3&12&13&85&58&$-73$
\end{tabular}
\end{center}

無限族は $t,k\geq1$ を整数として
\[
\begin{array}{c|c|c|c}
 (r,S,T,s,z)&(a,b)&(D_A,D_B)&\text{parameter}\\\hline
 (1,2,6,2t-1,t)&(4t-1,6t+1)&(3,-7)&t\geq1\\
 (2,3,6,2k,1)&(6k+1,7)&(6-6k,7k-9)&k\geq1\\
 (3,4,6,3k,2)&(12k+1,13)&(12-12k,13k-14)&k\geq1
\end{array}
\]
となる．第一族では二つの分母がともに増大する．
後の二族は固定分母七・十三の共鳴であり，公開稿の
\texttt{sp13:no-resonant-continuation} により，固定 ray をもう一つの二高分母 core に
接続して対向 $D$ を相殺することはできない．

\paragraph{導出：小さい reverse data から有限個の factorization へ．}
$h=S-r$，$k=T-r-1$ と置く．$h=0$ は canonical 補助変数 $\alpha$ の非零性に反する．
$h<0$ の可能性は $(r,S)=(3,2)$ だけである．
負の枝の二分により $T=S=2$ となり，整除性から $s=3u,z=2v$，$u,v\geq1$ と書ける．
このとき $\alpha=-u,\beta=-v$ であり，Bézout 式は
\[
 (9u+1)(-v)-(8v+1)(-u)=1,
 \quad\text{すなわち}\quad u(1-v)-v=1
\]
となるが，左辺は負である．従って六組
\[
 (r,S)=(1,2),(1,3),(1,4),(2,3),(2,4),(3,4)
\]
だけを調べればよい．全て $h,k>0$ である．

\[
 R=r(r+1),\quad U=(r+1)(S-r),\quad V=r(T-r-1),
 \quad E=r^2T-(r+1)^2S+R
\]
と定義すると，高高 pair product は
\begin{equation}\label{nr13:bilinear}
 Esz+Vz-Us=R
\end{equation}
と同値である．$s,z\geq1$ より
\begin{equation}\label{nr13:T-bound}
 s(E-U)\leq R-V
\end{equation}
である．一方，正の枝の不等式 $rT>(r+1)S$ を保持する．
これらから $T$ は次の二十六候補に限られる：
\[
\begin{array}{c|c}
 (r,S)&T\\\hline
 (1,2)&5,6,7\\
 (1,3)&7,8,9,10,11,12,13\\
 (1,4)&9,10,11,12,13,14,15,16,17,18,19\\
 (2,3)&5,6\\
 (2,4)&7,8\\
 (3,4)&6
\end{array}
\]
実際，$r=1$ で $T\geq6S-4$ なら $E-U\geq0$，$R-V=4-T<0$ となり矛盾する．
$(r,S)=(2,3)$ で $T\geq7$，$(2,4)$ で $T\geq9$ も同様である．
$(3,4)$ で $T\geq7$ なら
$s(E-U)\geq9T-56>24-3T=R-V$ となって矛盾する．

$E\ne0$ のとき，式~\eqref{nr13:bilinear} を
\begin{equation}\label{nr13:factorization}
 (Es+V)(Ez-U)=ER-UV
\end{equation}
と書く．$C=ER-UV\ne0$ なら $d=Es+V>0$ は $|C|$ の正約数であり，
\[
 s=\frac{d-V}{E},\qquad z=\frac{Us+R}{d}
\]
を調べればよい．正整数性，$r\mid sS$，$r+1\mid zT$，$a=sS+1,b=zT+1$ の奇数性を課すと，
表示した二十二模型がちょうど残る．候補生成と棄却理由は
\texttt{bounded\_reverse\_catalog\_verification\_v13.json} に全て記録している．

$E=0$ の三候補は $r=1,T=4S-2$ であり，式は
$(S-1)(2z-s)=1$ となる．従って $S=2,s=2z-1$ だけが残り，第一族を得る．
$E\ne0,C=0$ の二候補は $(r,S,T)=(2,3,6),(3,4,6)$ であり，
式~\eqref{nr13:factorization} と $Es+V>0$ から，それぞれ $z=1,z=2$ となる．
残る整除性が $s=2k,s=3k$ を与え，後の二族を得る．
全ての模型は元の pair products と row sums を満たし，
標準形の明示 weights により正である．

\paragraph{この分類の射程．}
有限表の最大分母八〇七は完全な factorization の結果として現れた数であり，
計算時の探索上限ではない．ただし reverse divisors の上限四は仮定である．
他の問題へ適用するときは，その上限が本当に node 条件等から従うことを先に示す必要がある．
また局所 charts の表は，整数 marked clutching，node 配分，actual areas の大域的実現を保証しない．

% v20: complete matching parameterization, with simultaneous nonnegative high D reduced to a nonsingular sector.
% Singular positive charts remain relevant to the broader local classification.
\subsection{三高分母の Schur 補行列と，符号原理における整数性}
\label{nsch18:three-high}

ここでは $q_0=1$，$q_1,q_2,q_3>1$ とし，
$A=\operatorname{diag}(q_i)+(p_{ij})$ が
\[
 A\mathbf1=0,\qquad q_iq_j-p_{ij}p_{ji}=1
\]
を満たす局所問題を扱う．目標は，整数模型において
$D_i=-q_i-\sum_{j\ne i}p_{ji}$，$i=1,2,3$，の少なくとも一つが負になるか，
少なくとも三つ全てが狭義正となることはないかを決めることである．
本項ではいずれも一般定理とは主張しない．

\paragraph{仮定を追加しない一般帰着．}
$r_i=p_{0i}$，$x_i=p_{i0}$ とすると
\[
 \sum_{i=1}^3r_i=-1,\qquad q_i=1+r_ix_i.
\]
全ての $r_i,x_i$ は非零であり，各対 $r_i,x_i$ は同符号である．
分母一の diagonal entry を pivot とした Schur 補行列
\[
 B=A_{\{1,2,3\},\{1,2,3\}}-xr^{\mathsf T}
\]
は整数行列で，
\begin{equation}\label{nsch18:unit-diagonal}
 B_{ii}=1,\qquad B\mathbf1=0.
\end{equation}
実際，diagonal は $q_i-r_ix_i=1$ であり，row sums は
$-x_i-x_i\sum r_j=0$ である．従って，ある整数 $\alpha,\beta,\gamma$ により
\[
 B=\begin{pmatrix}
 1&\alpha&-1-\alpha\\
 \beta&1&-1-\beta\\
 \gamma&-1-\gamma&1
 \end{pmatrix}
\]
と表せる．残る三つの pair equations は
\begin{equation}\label{nsch18:linear-x}
 x_i(r_i-B_{ji}r_j)+x_j(r_j-B_{ij}r_i)=B_{ij}B_{ji}
 \qquad(1\leq i<j\leq3)
\end{equation}
であり，$B,r$ を固定すれば $x$ に関する線形方程式である．
これは $p_{ij}=B_{ij}+x_ir_j$ を元の pair equation へ代入して得られる．
また，column sums を取ると
\begin{equation}\label{nsch18:defect}
 D_0=-1-\sum_i x_i,\qquad
 D_i=-\sum_{j=1}^3B_{ji}+r_iD_0.
\end{equation}
従って $D_0>0$ なら $\sum_{i=1}^3D_i=-D_0<0$ であり，
高分母側の少なくとも一つは直ちに負となる．
さらに，正の weights を持つ場合は全 $D$ ベクトルが零である整数模型は存在せず，
$D_0=0$ でも高側に負の係数が現れる．
従って正の weights の下で残る場合は $D_0<0$ に限られる（\ref{nhd20:reduction}）．

この行列問題は，二次の principal minors が全て一である非対称行列の特殊な場合である．
Lin--Sturmfels\cite[\S2, Proposition~4]{LinSturmfels} は principal minors と cycle sums の
整数多項式による変換を与える．本問題ではさらに row sums，整数性，正の weights を課すため，
一般の principal-minor relations だけから係数の符号や整除性は結論しない．

\paragraph{全ての正整数三高分母模型を尽くす canonical 表示．}
完全 matching 定理（\ref{nm18:matching}）により，正の numerators の辺は
高 rays 間の $ab$ と，残る高 ray $c$ と低 ray $0$ を結ぶ $c0$ に限られる．
従って以下の表示は追加の matching 仮定を必要とせず，全ての正整数三高分母模型を扱う．
これは有限個の模型への分類ではなく，残る整数条件を共通の線形系に表す帰着である．

$s,t,S,T,U\ge1$，$u=s+t-1$ と置く．適切に $a,b$ を交換すれば $k\ge0$ を選べ，
\[
 (q_a,q_b,q_c,q_0)=(1+sS,1+tT,1+uU,1),
\]
\[
 P=\begin{pmatrix}
 0&tS+\alpha&-(uS+\alpha+1)&-S\\
 sT+\beta&0&-(uT+\beta+1)&-T\\
 -(sU-k)&-(tU+k+1)&0&U\\
 -s&-t&u&0
 \end{pmatrix}
\]
となる．$H=tS+\alpha,K=sT+\beta,L=uS+\alpha+1,N=uT+\beta+1,
M=sU-k,O=tU+k+1$ の六つが全て正であることを保持する．
$\alpha,\beta$ は整数であり，
$\mathcal A=s\alpha-t,\mathcal B=t\beta-s,\mathcal C=s+uk$ と置く．
残る三つの pair equations は
\[
 \begin{pmatrix}
  \mathcal B&\mathcal A&0\\
  -\mathcal C&0&\mathcal A+1\\
  0&\mathcal C-1&\mathcal B+1
 \end{pmatrix}
 \begin{pmatrix}S\\T\\U\end{pmatrix}
 =\begin{pmatrix}
 -\alpha\beta\\k(\alpha+1)\\-(\beta+1)(k+1)
 \end{pmatrix}
\]
にちょうど同値である．正の weights は
$w=(q_a,H+1/(2K),u+1/(2U),1)$ と明示できる．
幾何的な奇数分母条件は $sS,tT,uU$ が全て偶数であることとして保持する．
公開算術稿の \texttt{sh14:matching-setup} と \texttt{sh14:linear-reduction} は，
この表示と線形化の証明箇所である．matching 定理により全正整数三高模型へ適用できる．

この canonical 表示では
\[
 \mathcal A=s\alpha-t\geq1,\qquad
 \mathcal B=t\beta-s\leq-2,\qquad
 \mathcal C=s+uk\geq2,\qquad \beta\leq-1.
\]
$d=U+1-S-T=-D_0$ と置くと
\[
 D_a=sd-\beta-k-1,\quad D_b=td-\alpha+k,\quad
 D_c=\alpha+\beta+1-ud.
\]
従って全ての高係数が非負となるには
\[
 sd\geq\beta+k+1,\quad td\geq\alpha-k,\quad
 ud\leq\alpha+\beta+1
\]
が必要十分である．全て狭義正なら右辺をそれぞれ一だけ強める．
境界 $\beta=-1$ では $k=0,D_0>0$ が証明され，その範囲から全高係数非負の模型は生じない．
全ての正整数三高分母模型のうち，高係数符号の問題で今後調べる範囲は
$b=-\beta\ge2$，$d\ge1$，$\alpha\ge k\ge b+1$，$\Delta>0$ である．
全ての特異系はこの符号問題から除かれる（\ref{nhd20:reduction}）．
matching 性が解決したことと，この $D$ 符号の問題が解決したことを混同しない．

固定 $s\geq2,t\geq1$ に対する $\beta=-1$ の全解は，
\[
 K_0=s^2+st+t^2-s-t,\quad g=\gcd(s-1,t),\quad G\mid K_0g,\quad G>0
\]
を列挙し，$\alpha=(G+K_0)/(st)$ の正整数性と
\[
 S=\frac{\alpha(s-1)(s\alpha-t+1)}G,\quad
 T=\frac{\alpha s(s+t-1)}G,\quad
 U=\frac{\alpha s(s-1)}G
\]
の正整数性を課して得る．幾何的な奇数条件は $sS,tT,uU$ が全て偶数であることとして最後に課す．
この完全な有限手続きは，固定 $s,t$ に関するものであり，一般分母の有限性は主張しない．

\paragraph{特異系は実在する：分類から除外しない．}
係数行列の determinant は
\[
 \Delta=-\mathcal B(\mathcal A+1)(\mathcal C-1)
          +\mathcal A\mathcal C(\mathcal B+1).
\]
$\Delta\ne0$ では Cramer 分子の整除性を確認する．$\Delta=0$ では拡大行列のrankと整数解格子を調べる．
実際，$(s,t,\alpha,\beta,k)=(3,12,12,-4,1)$ では rank が二で，正整数全解は
\[
 (S,T,U)=(136+200n,291+425n,93+136n),\qquad n\geq0
\]
で尽くされる．分母は $(409+600n,3493+5100n,1303+1904n,1)$，
高係数は $(-997-1467n,-4007-5868n,4671+6846n)$ である．
この固定sliceは解決済みであり，同じ補助データの探索を繰り返さない．
全ての特異系の分類と，大域の五 core 接着は未解決である．
ただし特異系の高係数が全て非負になる可能性は排除済みであり，その候補として再探索しない．
正整数三高分母の非matching型は存在しないことが証明されており，再探索の対象ではない．

\paragraph{高係数の同時非負性は非特異領域だけに残る．}
\label{nhd20:reduction}
canonical 表示で $b=-\beta$，$d=-D_0=U+1-S-T$ と置く．
整数 row sums，pair products，正の weights のみを仮定すると，
\[
 D_a,D_b,D_c\ge0
 \quad\Longrightarrow\quad
 b\ge2,\quad d\ge1,\quad \alpha\ge k\ge b+1,
\]
\[
 \Delta=\mathcal C(\mathcal A-\mathcal B)
          +\mathcal B(\mathcal A+1)
 \ge(-\mathcal B)(-\mathcal B-1)>0.
\]
分母の上限，奇数性，clutching，node 条件は不要である．
したがって，階数二を含む全ての特異 Schur 系から
全高 $D$ 非負の模型が生じる可能性はない．
この結論は，特異な正整数局所模型そのものを排除するものではない．

再利用する証明の要点を残す．
\[
 D_a=sd+b-k-1,\quad
 D_b=td-\alpha+k,\quad
 D_c=\alpha-b+1-ud,\qquad
 D_a+D_b+D_c=d
\]
である．まず全 $D$ 零を除く．その場合は
$d=0,\alpha=k=b-1$ で，$b\ge2$ である．
\[
 m=b^2-b+1,\qquad K=s^2+st+t^2-s-t,\qquad
 r=(-s,-t,u)^{\mathsf T},\quad x=(-S,-T,U)^{\mathsf T},
\]
\[
 B_0=\begin{pmatrix}
 1&b-1&-b\\
 -b&1&b-1\\
 b-1&-b&1
 \end{pmatrix}
\]
と置くと，Cramer 恒等式は $\Delta=mK$，$K=m-1$，
$x=r-B_0r/m$ を与える．整数性から各成分で $m\mid B_0r$ である一方，
\[
 B_0^{\mathsf T}B_0=m(3I-J),\qquad
 \lVert r\rVert^2=2m-1,\qquad \mathbf1^{\mathsf T}r=-1
\]
である．ここで $J$ は全成分が一の三次行列である．従って
$m^2\mid\lVert B_0r\rVert^2=m(6m-4)$，従って $m\mid4$ となる．
$m$ は三以上の奇整数だから矛盾する．
特に $D_0=0$ なら高係数には正と負がともに現れる．

同時非負性には $d\ge1$ が必要で，
$0\le D_a\le d$ から
\[
 (s-1)d+b-1\le k\le sd+b-1
\]
を得る．$k=b-1$ なら $s=1$，$u=t$ となり，
線形系の第三式は
\[
 t(b-1)T-tbU=b(b-1)>0
\]
となる．しかし $U=S+T+d-1>T$ なので左辺は負である．
$k=b$ なら $(s-1)d\le1$ であり，$s=1$ または $(s,d)=(2,1)$ となる．
前者では第三式 $tb(T-U)=b^2-1>0$ が同じ理由で矛盾する．
後者では $\alpha=b+t$，$U=S+T$ であり，第一式と第三式の線形結合が
\[
 t(1+b^2+tb)S=t+2b-b^3<t
\]
を要求するため，$S\ge1$ に反する．
この境界排除は determinant で割らないので，途中で特異系を落としていない．

以上から $k\ge b+1$ であり，
$D_c\ge0$ と $k\le sd+b-1$ から
$\alpha-k\ge(t-1)d\ge0$ が従う．
$\bar b=-\mathcal B=s+tb$ と置けば
\[
 \mathcal C-\bar b=(s-1)b+u(k-b)\ge0,\qquad
 \Delta=\bar b(\bar b-1)+(\mathcal C-\bar b)(\mathcal A+\bar b)
\]
により determinant の狭義正値性を得る．

残る非特異領域での一般符号問題は未解決である．
この領域の整数局所模型自体は存在する．例えば
\[
 (s,t,\alpha,\beta,k)=(1,2,10,-2,5),\qquad
 (S,T,U)=(4,5,11)
\]
は $q=(5,11,23,1)$，$D=(-1,1,3,-3)$，
正の weights $(5,109/6,45/22,1)$ を与える．
従って $\alpha\ge k\ge b+1$ と $\Delta>0$ だけで局所模型を捨てず，
Cramer 整除性と全高 $D$ の非負性を引き続き課す必要がある．
\nolinkurl{verify_high_d_boundary_v20.py} は
四十二の多項式恒等式と三つの厳密行列回帰を検算する．
一般命題の根拠は上の恒等式・不等式であり，有限範囲の代入結果ではない．

\paragraph{五つの整数を固定した三高 $D$ の完全復元．}
\label{nhd21:reconstruction}
同時非負性の残余問題では，$x=D_a,y=D_b,z=D_c\ge0$ を直接指定する．
$s,t$ も固定することが必要であり，$D$ だけを固定した有限性とは主張しない．
次の変数は再開後の復元と消去にそのまま使える：
\[
 u=s+t-1,\quad d=x+y+z\ge1,\quad
 k=sd+b-1-x,\quad\alpha=ud+b-1+z,
\]
\[
 h=s-x,\quad j=t-y,\quad\ell=u+z,\quad
 \lambda=ux+sz,\quad\mu=uy+tz,
\]
\[
 F=b-1+\lambda,\quad m=F^2+F+1,\quad
 E=\mu h-\ell-j,\quad J=h+\ell-\lambda j.
\]
正の模型では $h\ge1$，すなわち $D_a\le s-1$ が新しく必要となる．
$j$ の符号は制限しない．$E,J$ は $b$ に依存せず，
\[
 S=h+\frac{FE-J}{m},\quad
 T=j+\frac{E+(F+1)J}{m},\quad
 U=\ell+\frac{(F+1)E+FJ}{m}
\]
が全尺度を一意に復元する．三尺度の整数性は単一条件 $m\mid FE-J$ と同値である．
元の三つの高高 pair 残差はこの尺度で互いに等しくなるが，自動的に零にはならない．
\[
 (-s-tb)S+(s\alpha-t)T-\alpha b=0
\]
を最後まで保持し，全尺度の正値性，canonical 六符号量の正値性，
$\alpha\ge k\ge b+1$ も確認する．これらは元の模型への必要十分判定となり，
正の weights は公開稿の明示式から得られる．

\paragraph{ノルムによる全 $b$ の有限化．}
$v=S-h,w=T-j$ と置くと
\[
 E=(F+1)v+w,\qquad J=-v+Fw,
\]
\[
 N=E^2+EJ+J^2=m(v^2+vw+w^2).
\]
零ノルムは，正尺度の下では
$(s,t,x,y,z)=(3,3,0,0,1),(2,2,1,1,1)$ に限られる．
元の残余式がそれぞれ $b(b+5)=27$，$b(b+9)=14$ を要求するため，両方とも排除される．
従って $N>0$ であり，
\[
 m\mid N,\qquad m\le N,\qquad
 b\le 1-\lambda+
 \left\lfloor\frac{\sqrt{4N-3}-1}{2}\right\rfloor
\]
となる．実装では平方根を \texttt{isqrt} で厳密に評価する．
$s,t,x,y,z$ を固定すれば，任意 cutoff なしに $b\ge2$ を全て調べられる．
この五つの整数そのものへの一様上界は未証明である．

\paragraph{誤った簡約を防ぐ二つの反例．}
$(s,t,x,y,z,b)=(1,1,0,3,0,3)$ では
$F=2,m=7,E=4,J=2,N=28$ である．$m\mid N$ でも
$m\nmid FE-J=6$ なので尺度は整数でない．ノルム整除性だけに置き換えない．
また $(1,3,0,4,0,2)$ では $(S,T,U)=(3,5,11)$ は正整数で，
六符号量も正，$D=(0,4,0,-4)$ であるが，元の三つの高高 pair 残差は全て三である．
この例は模型ではなく，残余式を落とせないことの確認例である．

\paragraph{再利用する実装と次の課題．}
\nolinkurl{verification/verify_high_d_reconstruction_v21.py} の
\texttt{reconstruct\_fixed\_data(s,t,x,y,z)} が全 $b$ の完全復元器である．
通常の実行は二十四の多項式係数恒等式と五つの厳密回帰を確認する．
\texttt{--fixed 3 4 0 1 1} のように固定データを渡せる．
既存の有限族検査は $s,t\le12,0\le x<s,0\le y,z\le4,d\ge1$ の
$23256$ 組について $b\ge2$ を全範囲で尽くし，模型は零だった．
結果と実行当時の source は \nolinkurl{research/} に保存した．
これは五変数に制限を付けた結果であり，一般不存在ではない．
第22版ではこの復元式に entry 正値性を加え，固定するデータを $(s,t)$ の二つへ減らす．
次の節の有限化と合わせて用い，二変数自体が動く問題を残す．

% v22 research notes: retain only material useful for future work.
\subsection{三高同時非負性：低行二変数までの有限化}
\label{notes:hd22}

\paragraph{現在の到達点}

第21版の五固定データ $(s,t,x,y,z)$ ごとの $b$ 有限化は，
第22版で固定 $(s,t)$ ごとの全 $x,y,z,b$ の有限復元へ強まった．
新しい主モジュールは
\nolinkurl{high_d_fixed_low_row_finiteness_v22.tex} である．
低行 $(-s,-t,s+t-1,0)$ の二整数自体の一様上限，
および三高同時非負模型の一般不存在は未解決である．

利用する記号と核となる式だけを残す．
\[
 u=s+t-1,\ h=s-x\ge1,\ B=b-1\ge1,\
 F=B+ux+sz,\ m=F^2+F+1.
\]
固定 $(s,t,x,z,b)$ において
\[
 a=uh-B(F+1),\qquad
 n_0=(F+1)(Bt+h)+F(u+z)+htz,
 \qquad mT=n_0+ay.
\]
$K=sT-b>0$ と $y\ge0$ が
\[
 \boxed{F\le M:=st+(2u+1)h-2}
\]
を与える．上限の証明は $T$ の正整数性を過剰に使わず，entry $K$ の正値性を
直接使う．$F$ は整数なので，$F\ge M+1$ の排除から表示した上限が従う．
さらに $b-1+sz\le M-ux$ より
\[
 (3u+1)x\le st+(2u+1)s-3
\]
も従う．これは $x\le s-1$ の追加の絞り込みとして使える．

残余式を $y$ の多項式にして
\[
 \mathcal R(Y)=suaY^2+(aL-C)Y+C_0
\]
とすると，$L,C>0$ であり，$a=0$ でも一次係数が $-C<0$ となる．
従って各 $(s,t,x,z,b)$ に非負整数根は高々二個，$a\le0$ なら高々一個である．
二次式の係数と $C_0$ の計算法は公開用証明に明示した．
根は判別式の整数平方根と正確な割り算で全て求められ，$y$ の探索上限は不要である．

\paragraph{完全復元で残す検査}

二次式の根であることは元の第一残余そのものの零を意味する．
ただしそれだけでは整数模型ではない．必ず
\[
 d=x+y+z\ge1,\quad \alpha\ge k\ge b+1,\quad m\mid FE-J,
\]
三尺度の正値性，六つの符号量の正値性を検査する．
最後に元の $q,P$ を組み立て，row sums，六つの pair products，
$D=(x,y,z,-d)$，正 weights に対する十二の gaps を再計算すれば，
導出式の取り違えも検出できる．$j=t-y$ の非負性は仮定しない．

第21版の零ノルム分類は新手順の例外処理として不要になったが，
$(s,t,x,y,z)=(3,3,0,0,1)$ と $(2,2,1,1,1)$ の両族は
元残余を捨ててはいけないことを示す回帰例として残した．
前者の尺度 $(3,3,6)$ は $b$ によらず整数だが，
残余 $27-b(b+5)$ は $b\ge2$ の整数で零にならない．

\paragraph{再現可能な厳密計算と次の分岐}

標準ライブラリのみの
\nolinkurl{verify_high_d_low_row_finiteness_v22.py}
は二十四の多項式恒等式を係数ごとに確認する．
\texttt{--fixed s t} はその低行の全模型を完全復元し，
\texttt{--scan 12} は $1\le s,t\le12$ の全 $144$ 低行に同じ処理を行う．
探索箱がかかるのは $s,t$ だけであり，$D$ や分母は全範囲である．
結果は \nolinkurl{high_d_low_row_finiteness_verification_v22.json} にある．
全五組候補 $2{,}004{,}015$，$a=0$ の一次式 $219$，
非負整数残余根六個で，全て $d=0$ のため対象外だった．
零 $D$ の六根の $(s,t,b)$ は
\[
 (3,4,6),(4,3,6),(6,9,13),(9,6,13),(10,12,19),(12,10,19).
\]
これらの尺度は有理数であり，整数の零 $D$ 模型の反例ではない．

次の有力課題は，無上界の $(s,t)$ に対して残余二次式の非負整数根と
単一合同条件を同時に扱うことである．有限復元を使えば，
単純な分母・高係数 cutoff の拡大を避け，低行の構造に沿った反例探索ができる．
また観測範囲では $d\ge1$ の整数残余根が合同条件を課す前から存在しなかった．
これは，尺度の整数性を外した有理復元段階にもより強い障害がある可能性を示すが，
現時点では観測であって定理ではない．
判別式は
\[
 \mathcal D_y=(aL-C)^2-4suaC_0
 =C^2+a\bigl(aL^2-2LC-4suC_0\bigr)
\]
である．一般の全展開では扱いやすい積を得ていない．
次は $a$ の符号と $a=0$ の退化面を先に分け，
整数平方間の評価と合同条件を組み合わせる方がよい．
$a=0$ は $B(F+1)=uh$，すなわち
\[
 u(s-bx)=B(b+sz)
\]
を要求するので，$s-bx>0$ という追加制約も得られる．
この面で残余は $-CY+C_0$ であり，$y=C_0/C$ の整数性へ直ちに帰着する．

判別式とは別に，$s=t$ かつ $x=y$ の対称 sector は無上界で排除できた．
このとき $J=-E$，単一合同条件から $n=E/m\in\mathbb Z$，
$S=h+(F+1)n$，$T=h-Fn$ となる．
$F\ge h+1$ により正値性は $n=0$ を強制し，既知の零ノルム排除に矛盾する．
この証明は整数性と正値性の両方を使うので，
一般 sector でも $E+J=\ell(sy-tx)+h-j$ の小さな値から類似の格子拘束を
引き出せるかを検討する価値がある．
一般不存在を目標に固定せず，低行に沿った無限族・合同条件の反例・
残余式の因数分解が現れた場合は，その構造の記述へ研究計画を切り替える．

% v23 notes: keep only reusable formulas, exact scope, and useful next steps.
\subsection{三高同時非負性：非対称量と退化面}
\label{notes:hd23}

\paragraph{確定した四領域と残る量化}

主証明は \nolinkurl{high_d_asymmetry_lattice_obstruction_v23.tex} にある．
第22版の固定低行 $(s,t)$ ごとの全係数有限復元を保ち，
\[
 \delta=sy-tx,\qquad a=uh-(b-1)(F+1)
\]
を使って，次の四領域を新たに無上界で除外した：
\[
 \delta=0,\ s\ge t;\qquad
 \delta=-1;\qquad
 \delta=1,\ x\ge1;\qquad
 a=0,\ s\le2.
\]
$s=t$ だけで模型を除外したわけではない．
$\delta=0,s<t$，$|\delta|\ge2$，$\delta=1,x=0$，
一般の $a=0,s\ge3$ は残る．三高同時非負性の一般不存在は未解決である．

再利用する核は，$n=\ell-U\in\mathbb Z$ に対する
\[
 S=j-(F+1)n-\ell\delta,\qquad
 T=h+Fn+\ell\delta,
 \qquad F+\delta=b-1+h+s(d-1)
\]
である．$n$ は node 配分ではなく尺度の格子変数である．
$\delta=-1$ では正値性が $x=1,n=1,t=sy+1$ を強制し，
残る差方程式が正係数表示を持つ．
$\delta=1,x\ge1$ では $x=1,n=-1,z=0,t=sy-1$ を強制する．
後者では差方程式 $Q_+=0$ に加えて元残余を残すことにより
\[
 y\mathfrak r=\bigl(sy(y+1)+1\bigr)Q_++N_+,
 \qquad N_+>0
\]
を得る．この二方向の違いは今後の $|\delta|\ge2$ でも保持すべきである．

\paragraph{差方程式だけでは通過する新しい回帰例}

\[
 (s,t,x,y,z,b)=(5,4,1,1,0,3)
\]
は $\delta=1$ であり，復元尺度は $(S,T,U)=(6,2,9)$，
$\alpha=18,k=11$，$q=(31,9,73,1)$ となる．
単一合同と六つの符号条件を満たすが，三つの元の高--高 pair 残余は全て $16$ である．
これは模型ではない．単一合同と正値性に元残余を付け加える必要を具体的に示す．
また $(s,t,x,y,z,b)=(3,6,1,2,0,3)$ では
$\delta=0,s<t$，$(S,T,U)=(4,2,8)$，元残余 $-18$ である．
こちらも模型ではなく，今回の比例領域の順序仮定を勝手に捨てる根拠にならない．

\paragraph{退化面の有界整数と追加の整除条件}

$a=0$ では
\[
 s\ell=bF,\qquad uh=(b-1)(F+1),\qquad
 u(s-bx)=(b-1)(b+sz)>0
\]
となる．公開証明の補題により
\[
 K_0=h(s-1)+1,\qquad R_0=b^2-sT
\]
は
\begin{equation}\label{notes:hd23-degenerate-quadratic}
 \begin{gathered}
 1\le R_0\le K_0,\\
 R_0F^2+(R_0-bK_0+sh)F+R_0-b^2+sh=0
 \end{gathered}
\end{equation}
を満たす．$s=1,2$ では $R_0$ はそれぞれ一通りと三通りだけとなり，
各判別式の法三，五，七の平方障害で全て除外できた．

元残余も使うと，さらに
\begin{equation}\label{notes:hd23-extra-divisibility}
 \boxed{\ b-1\mid h(s-1)(R_0-1)\ }
\end{equation}
を得る．この式の短い導出を残す．$B=b-1$ とし，
\[
 S_0=FT-Bt-\ell,\quad \alpha_0=F+tz,\quad
 H=s+tb,\quad K=sT-b,\quad \rho=BH-uK
\]
と置く．$S=S_0+By$，$\alpha=\alpha_0+uy$ だから
\[
 \mathfrak r=\mathfrak r_0-\rho y,
 \quad
 \mathfrak r_0=-HS_0+(s\alpha_0-t)T-\alpha_0b.
\]
$uh=B(F+1)$ より $B\mid h\rho$ なので，$\mathfrak r=0$ なら
$B\mid h\mathfrak r_0$ である．法 $B$ で $b=1$ とすると，恒等式
\[
 h(\mathfrak r_0+tT)
 =h(tT+1)(s\ell-F)+t(1-sT)uh
 \pmod B
\]
の右辺は零である．従って $B\mid htT$ であり，
$ht\equiv h(1-s)\pmod B$ から $B\mid h(s-1)T$ となる．
最後に $sT=b^2-R_0\equiv1-R_0\pmod B$ を使えば
式~\eqref{notes:hd23-extra-divisibility} が従う．

これにより，退化面の大部分は低行の一整数 $s$ だけを固定して有限化できる．
$x\ge1$ の場合は $s-bx\ge1$ から
\[
 2\le b\le\left\lfloor\frac{s-1}{x}\right\rfloor.
\]
$x=0$ かつ $R_0\ge2$ では $h=s$ なので
\[
 2\le b\le1+s(s-1)(R_0-1)\le1+s^2(s-1)^2.
\]
$s,h,b,R_0$ を固定すれば，式~\eqref{notes:hd23-degenerate-quadratic} の
正の二次係数 $R_0$ により整数 $F$ は高々二個であり，
\[
 u=\frac{B(F+1)}h,\qquad \ell=\frac{bF}s,\qquad
 t=u-s+1,\quad x=s-h,\quad z=\ell-u
\]
で残りを復元できる．$y$ は第22版の恒等零でない一次残余で決まる．
全ての整数性，正値性，単一合同，元の pair 残余は引き続き確認する．

この有限化から外れる唯一の組合せは $x=0,R_0=1$ である．
この枝には必要条件
\begin{equation}\label{notes:hd23-exceptional-congruence}
 s^2\mid2(b-1)
\end{equation}
がある．実際，$s\mid b(b-1)$ と $sT=b^2-1$ を比較すると $s\mid b-1$ であり，
$b-1=s\kappa$，$F=sL$ と書ける．
式~\eqref{notes:hd23-degenerate-quadratic} は
\[
 s\bigl(L^2+(1-\kappa K_0)L-\kappa^2+1\bigr)=2\kappa
\]
を与える．従って $s\mid2\kappa$，すなわち式~\eqref{notes:hd23-exceptional-congruence} が従う．
この枝の一般不存在は証明していない．平方判別式と元残余を併用する，
より小さい次の標的として扱える．

$a=0$ だから全ての整数尺度が存在しない，という強化は誤りである．
\[
 (s,t,x,y,z,b)=(3,138,0,0,4,16)
\]
では $F=27,m=757$，$(S,T,U)=(54,84,141)$，$R_0=4$ となり，
単一合同と六符号は成立する．元残余は $5658$ であり模型ではない．
この例は一般の退化面で元残余を使う必要性を示す．

\paragraph{残る $\delta=1,x=0$ の二整数による完全復元}

この枝では $s=y=1$ である．$k=b+z\ge b+1$ より $z\ge1$ であり，
格子正値性から $p=-n\ge1$ と書ける．
\[
 F=b-1+z,\quad m=F^2+F+1,\quad
 S=(F+1)p-z-1,\quad T=t+z+1-Fp,\quad U=t+z+p.
\]
差方程式は
\[
 (t-1)(z-1)=mp-F(z+2)
\]
である．$z=1$ なら $mp=3F$ だが，$F=b\ge2$ で
$m-3F=(F-1)^2>0$ なので不可能である．従って未解決枝は $z\ge2$ に限られる．

固定 $b,z\ge2$ に対して $D=z-1$ とし，整数 $p$ の一次式を
\[
 \tau=D+mp-F(z+2),\quad
 \sigma=(F+1)p-z-1,\quad
 \theta=\tau+D(z+1-Fp)
\]
と定める．$t=\tau/D$，$S=\sigma$，$T=\theta/D$ であり，
元残余の分母を払った式は
\begin{equation}\label{notes:hd23-axis-quadratic}
 \begin{aligned}
 \mathcal P(p)={}&-D(D+b\tau)\sigma
 +(z\tau+FD)\theta\\
 &-bD\bigl((z+1)\tau+FD\bigr)=0.
 \end{aligned}
\end{equation}
その二次係数は
\[
 (b^2+bz+z^2)m>0
\]
なので，各固定 $(b,z)$ に正整数根 $p$ は高々二個である．
整数平方根による全根計算の後，
\[
 D\mid\tau\quad\Longleftrightarrow\quad
 z-1\mid(b^2+b+1)p-3b
\]
と $t,S,T,U$ および六符号の正値性を検査すればよい．
ここでは差方程式と元残余の双方を保持しており，通過したものは
第21版の完全復元により全 pair products と正 weights を持つ．
尺度整数性の単一合同も直接再確認できる．

標準 Python 検証器の \texttt{remaining\_axis\_fixed\_b\_z(b,z)} はこの完全復元を実装する．
\texttt{--axis b z} でも呼び出せる．$p$ や $t$ に探索上限はない．
固定している $b,z$ 自体には上限を与えていないので，この手順だけで残余枝全体を
不存在と結論してはいけない．

\paragraph{検算と今後の焦点}

\nolinkurl{verify_high_d_asymmetry_v23.py} は標準ライブラリのみを用い，
四十一の多項式恒等式を係数ごとに検査する．
退化面の四つの判別式は，法 $9,25,9,49$ の全九十二剰余類で
平方剰余と交わらないことを独立確認する．七組の回帰例は
元の rows，pair products，$D$，必要な場合の正 weights を再計算する．
結果は \nolinkurl{high_d_asymmetry_verification_v23.json} にあり，
source と公開証明の hash，検算の量化，残る範囲を記録している．
第22版の二百万余の復元計算は再実行していない．

次は，$\delta$ の絶対値を単に切って列挙するのでなく，
格子正値性が $x,n$ をどう制限するかを一般の $\delta$ で整理する．
とくに負側 $\delta=-k$，$s\ge k$ では，$j>0,n\ge1$ と $S>0$ から
\[
 u(x-k)+(s-k)z<t-(b-1)-1
\]
を得るので，$x\le k$ が必要である．
正側でも $\delta=k$，$s\ge k$ なら $n\le-1$ と $T>0$ が $x\le k$ を与える．
この一般の狭い格子幅と，二次残余の符号証明を組み合わせることが有望である．
また，退化面の $x=0,R_0=1$ と，正側零境界の式~\eqref{notes:hd23-axis-quadratic} は
独立した二つの具体的な未解決問題として追える．
どちらも反例や無限族が現れれば，その構造の記述へ計画を変更する．


\paragraph{整数性を外すと，三つの高 $D$ が全て狭義正になる．}
次は\emph{有理数による代数的緩和}の反例であり，整数 branch denominator を持つ
局所 chart の反例ではない．頂点順を $(a,b,c,0)$ とし，
\[
 q=\left(\frac{79}{25},\frac{128}{25},\frac{274}{25},1\right),\qquad
 P=\begin{pmatrix}
 0&179/25&-231/25&-27/25\\
 53/25&0&-259/50&-103/50\\
 -91/25&-266/25&0&83/25\\
 -2&-2&3&0
 \end{pmatrix}
\]
と取る．全 row sums と pair equations を満たし，
\[
 w=(1,7/3,13/4,1)
\]
について全ての $w_j-w_ip_{ij}/q_i$ は狭義正である．
実際，正の辺は $a\leftrightarrow b$ と $c\leftrightarrow0$ のみで，
対応する四つの差は
\[
 \frac{16}{237},\quad\frac{13}{384},\quad
 \frac{17}{1096},\quad\frac14
\]
である．一方，
\[
 D=\left(\frac9{25},\frac9{25},\frac{23}{50},-\frac{59}{50}\right).
\]
この例でも low row と Schur 補行列
\[
 B=\begin{pmatrix}1&5&-6\\-2&1&1\\3&-4&1\end{pmatrix}
\]
は整数である．従って，$B$ と low row の整数性まで仮定しても，
$q$ と逆向き numerators の整数性を外すと一般符号原理は成立しない．
正の weights と主小行列の不等式だけによる証明は不可能であり，
残る整除性を本質的に用いる必要がある．
この反例と上記恒等式は \texttt{verify\_three\_high\_schur\_v12.py} で有理数演算により再確認できる．

\paragraph{$\beta=-1$ を持つ無限族と，一つの slice の完全な記述．}
\label{nsch18:beta-minus-one}
上の完全約数表示の $t=1$ の場合には，さらに簡単な閉じた表示がある．
一般のmatching型に $\beta\leq-2$ を仮定してはいけないことも，この族から分かる．
実際，$t=1$，$\beta=-1$，$k=0$ と固定した完全マッチング型は，
次の二つの正整数 $d,e$，$de\geq2$ でちょうど記述される：
\begin{align*}
 s&=de,&\alpha&=d(e+1),\\
 S&=d(e+1)^2(de-1),&
 T&=de(e+1),&
 U&=(e+1)(de-1).
\end{align*}
分母と numerators は
\[
 (q_a,q_b,q_c,q_0)=(1+sS,1+T,1+sU,1),
 \qquad
 P=\begin{pmatrix}
 0&S+\alpha&-(sS+\alpha+1)&-S\\
 sT-1&0&-sT&-T\\
 -sU&-(U+1)&0&U\\
 -s&-1&s&0
 \end{pmatrix}
\]
である．三つの高分母は全て奇数で，一より大きい．
$H=S+\alpha$，$K=sT-1$ と書くと，正の weights は例えば
\[
 w=\left(q_a,\ H+\frac1{2K},\ s+\frac1{2U},\ 1\right)
\]
で与えられる．正の四つの directed edges に対応する basin gaps は
\[
 \frac1{2K},\quad\frac1{2q_b},\quad\frac1{2U},\quad\frac1{2q_c}
\]
であり，残る edges は負の numerators を持つので正値性は自明である．
また
\[
 D_0=S+T-U-1=(\alpha-1)U+T-1>0,
 \qquad
 (D_a,D_b,D_c)=(-sD_0,-D_0-\alpha,sD_0+\alpha).
\]
従って，これは三つの高係数が全て正となる反例ではない．

完全性の確認も短い．$t=1,\beta=-1,k=0$ の線形系では
\[
 \Delta=s(s-\alpha),\qquad
 S=\frac{\alpha^2(s-1)}{\alpha-s},\quad
 T=\frac{\alpha s}{\alpha-s},\quad
 U=\frac{\alpha(s-1)}{\alpha-s}.
\]
$s\geq2$ であり，$\alpha=s$ の場合には $S$ の Cramer numerator
$-\alpha^2s(s-1)$ が非零なので系は不整合である．
$S>0$ には $d=\alpha-s>0$ が必要である．
$T,U$ の整数性は $d\mid s^2$ および $d\mid s(s-1)$ と同値なので，$d\mid s$ が従う．
逆に $d\mid s$ なら三つ全てが整数となる．$s=de$ と書けば表示した族を得る．
pair equations と row sums はそのまま代入して確認できる．

小さい具体例として $d=1,e=2$ を取ると
\[
 q=(19,7,7,1),\qquad
 P=\begin{pmatrix}
 0&12&-22&-9\\
 11&0&-12&-6\\
 -6&-4&0&3\\
 -2&-1&2&0
 \end{pmatrix},\qquad
 D=(-22,-14,25,11),
\]
かつ $w=(19,265/22,13/6,1)$ となる．
これは既存の有限 catalog にも含まれる模型であり，この slice 全体の無限パラメータ表示と，残る整除性の特徴づけを再利用できる．
これらの恒等式と具体例は \texttt{verify\_three\_high\_bounds\_v13.py} で厳密演算により検算する．

\subsection{次の接着計算に使う二つの局所障害}
\label{ng16:local-obstructions}

二価 core $i$ の二つの高近傍を $j,k$ とする．$q_{jk}=1$ なら，
core $i$ の二高分母標準形は必ず正の枝である．従って
\[
 S>r,\qquad rT>(r+1)S,\qquad L=T-S+1\geq3.
\]
この結論に対向 $D$ の相殺は不要であり，core $j,k$ の高分母次数を制限しない．
負の枝 $S=T=w<r$ を用いるときは，その core を含む高分母三角形が必要である．
一般の三高分母中心に二高分母の真約数補題を適用する代わりに，
中心側では signed divisors と高辺の合同式を直接比較する．
負の枝の $H=r-w$，$c=s-x$，$e=z-y$ に対して
\[
 Ece=e-c+1,\qquad E=1-\frac{w^2}{H(H+1)},
 \qquad -1<E<\tfrac12,\quad E\ne0
\]
となり，中心間が低辺なら reverse 和が $E-1$ となることが主な障害である．
例外 $e=1$ も node 数 $2+rw/H>6$ により排除される．

全低分母 core $*$ がある場合には，残る四頂点を $I$ とし
\[
 x_{ij}=p^i_{j*},\qquad r_{ij}=p^i_{*j}
\]
と置く．四辺 $i*$ の $D$ 相殺，局所の行和，および高辺の整数 clutching から
\[
 \sum_{i\ne j\in I}x_{ij}=\sum_{i\ne j\in I}r_{ij}=-4,
 \qquad
 \sum_{\substack{i<j\in I\\q_{ij}=1}}
 (x_{ij}+x_{ji}-r_{ij}-r_{ji})=0
\]
を得る．全低 core の自由な entries を固定しなくても使える総和条件である．
diamond では和に低辺が一つだけ残り，二価 core の低低 entries が
反対向き・同じ $L\in\{1,3,4\}$ にそろう．三つの場合は全て排除済みである．
$K_4$ 型ではこの和は空になるが，四本の $D$ 相殺を頂点別に保持すると，
後述する通常辺上の流れの式を得る．


\subsection{共通低 puncture における通常接着の上界}
\label{ng16:ordinary-bound}

$q_{ij}>1$，$q_{ik}=q_{jk}=1$ とし，高辺 $ij$ の puncture $k$ が通常の反対符号接着なら
\[
 p^i_{jk}=-p^j_{ik},\qquad p^i_{kj}=-p^j_{ki}=u,\qquad1\le |u|\le4.
\]
core $k$ の低低 entries を $a=p^k_{ij},b=p^k_{ji}$ とすれば $ab=0$ であり，
二つの node 数 $n,m\in\{0,\ldots,5\}$ に対して
\[
 a=1-n-u,\qquad b=1-m+u,\qquad
 0=-\{u-(1-n)\}\{u-(m-1)\}
\]
となるためである．高分母次数と対向 $D$ の相殺には依存しない．
node 総数が同じ $N\ge2$ の模型では，同じ証明で $|u|\le N-1$ を得る．
最大・単位約数の同符号例外をこの上界に含めてはいけない．

\subsection{一高分母の自由度を保持する標準形}
\label{ng16:leaf-data}

一高分母 core は，高 ray $A$ と低 rays $B,C,D$ の適切な標識により
\[
\begin{gathered}
 m=th(h+1),\quad q=m+1,\quad u=th,\quad v=t(h+1),
 \quad h,t\geq1,\\
 P=\begin{pmatrix}
 0&-m&-h-1&h\\
 -1&0&0&0\\
 -u&u-1-c&0&c\\
 v&-v-1-d&d&0
 \end{pmatrix},\qquad c,d\in\ZZ,\quad cd=0
\end{gathered}
\]
と書け，全ての整数 choices に正の weights が存在する．係数は
\[
 (D_A,D_B,D_C,D_D)=(-m-t,\ m+t+1+c+d,\ h-d,\ -h-1-c).
\]
$q=3$ では二つの負の reverse entries が一致し得る．
高 row だけを固定して低低 entries を一つの模型に置き換えず，
この $c,d$ の自由度と必要な低 ray の交換を保持する．

二高分母の低係数は
\[
 D_C=y-x+T-S,\qquad D_D=s-z-1
\]
であり，同時に負とはならない．一高分母の負 reverse を持つ ray $E$ は $B$ または $C$ である．
別の低 ray $J$ から $p_{JE}\leq-1$ を要求される場合，
$E=C$ なら $D_E>0$，残る低 ray から $E$ への entry は零となる．
$E=B$ かつ $D_E\leq0$ なら，残る低 ray からの entry は
$t(h+1)^2-D_E\geq4$ である．
係数相殺と node 条件を組み合わせるときに使える二分である．

\subsection{二葉付き三角形の完全な有限化から再利用する式}
\label{ng16:two-leaf-certificate}

$H=\{01,02,12,03,14\}$ は全排除済みである．再利用するのは，
未知の中心分母を走査せず，葉の自由度を保存したまま有限化する次の手順である．
中央 core $2$ の交差標識 $(A,B,C,D)=(0,1,4,3)$ と正の枝に対して，
共通低 puncture の上界は
\[
 1\le r\le4,\qquad z=1
\]
を同時に要求する．$T\ge5$ により辺 $12$ は正の最大・単位約数の例外となり，
\[
 b=T+1,\quad p^1_{23}=T,\quad p^1_{32}=1,\quad
 p^3_{12}=0,\quad n_{12;3}=n_{13;2}=0
\]
を得る．$r>1$ なら，辺 $02$ には通常接着が強制される．

$T=(r+1)y$，$E=D^2_4$ と置くと，
\begin{align*}
 (r+2)sS-r(rs+1)y&=2r(s-1),\\
 F_r(s)S&=r(rs+1)E+2r(r+2)(s-1),\\
 F_r(s)&=-rs^2+(4r+3)s-r,\\
 E&=(r+2)y-\frac{s+r}{r}S.
\end{align*}
$r=1,\ldots,4$ では $F_r(s)\ne0$ であり，$E\ge-2$ なら $s\le9$ となる．
より一般に，整数の正枝で $z=1$ なら，任意の正整数 $r,s$ でも $F_r(s)\ne0$ である．
実際，$F_r(s)=0$ は $r=3s/(s^2-4s+1)$ を与え，正整数解は $(r,s)=(12,4)$ だけだが，
上の係数式は $E=-12/7$ を要求して整数性に反する．
従って $(r,s,E)$ は $S$，続いて $y,T$ を一意に復元する．
また $E\ge-M$，$M\in\ZZ_{\ge0}$ なら
\[
 E<\left(4+\frac4r-s\right)\frac Sr,\qquad S>r
\]
から $s\le M+7$ を得る．これは他の係数下界が得られた場合にも再利用できる．
$S\ge4$ の下で，通常の $C/D$ 表現を持つ葉 $4$ では，
$-2\le E\le7$ の枝と，$r=1$，$D^4_3=E+\kappa$，
$\kappa\in\{2,5,8,10\}$ の枝に分かれる．
後者でも葉 $3$ の一方の低標識なら $th\le4$，
他方なら $h\le12,t\le6$ が証明される．
小さい $S$ の共鳴族 $(r,S,T)=(2,3,6)$ では，$s$ は正の偶数であり，
$s\ge8$ は他葉に向かう二つの正の $D$ が相殺しないため排除される．

この有限化は三十一中央模型・六十六生成 witness を与え，
全葉を復元すると十中央模型・五十二低三角形だけが残る．
葉 $4$ の ray $2=B$ を含む別の有限化は，
三十三中央模型・六十二 witness から四中央模型・七組を与える．
この七組は五十二組の部分集合であり，個数を加えてはいけない．
分母三の同じ標識付き行列は，一度だけ数える．
$B$ 表現の例外を省くことも，低低自由度 $c,d$ を固定することも許されない．

\paragraph{既知三行だけによる中心の延長障害．}
三つの隣接 cores の行列が既知なら，中心 $i$ の row $j$ は
\[
 p^i_{jk}=q_{ij}(1-n_{ij;k})-p^j_{ik}
\]
で復元される．非負 node 数の三分割は二十一通りなので，既知三行は全 $21^3$ 通りである．
三つの既知 rays の pair equations は
\[
 p^i_{jk}p^i_{kj}=q_{ij}q_{ik}-1
\]
である．五十二組のうち五十一組は中心 $0$，残る一組は中心 $1$ において，
この三行条件だけで候補が零となる．
未知の中心間分母，第四行，中心での $D$ 相殺，正の weights の判定は，
この最後の矛盾に必要ない．
この手順は，他の残存 graph で複数の隣接行が既知となる場合にも使える．

\subsection{\texorpdfstring{$K_4$}{K4} の低 hub と通常辺上の流れ}
\label{ng16:k4-flow}

全低 core を $4$ とし，$I=\{0,1,2,3\}$ の高 $K_4$ で
\[
 R_{ij}=p^i_{4j},\quad X_{ij}=p^i_{j4},\quad V_{ij}=p^4_{ij},
 \quad N_{ij}=1-R_{ij}-V_{ij}
\]
と置く．各行の node 予算は $\sum_jN_{ij}=5$ である．
通常辺では $R_{ji}=-R_{ij}$，$1\le|R_{ij}|\le4$，
例外辺では $R$ の対が $\pm\{1,q-1\}$ となる．
hub に接する $D$ 相殺を頂点別に保持すると
\[
 \sum_{j\ne i}(X_{ij}-R_{ji})
 =\sum_{j\ne i}N_{ji}-5
\]
を得る．例外辺の寄与は零であり，通常辺の寄与は
\[
 \operatorname{sgn}(R_{ij})W_{ij},\qquad
 W_{ij}=\frac{q_{ij}-1}{|R_{ij}|}+|R_{ij}|
\]
である．proper cut の発散の絶対値は高々十なので，通常 graph の橋では
$W\le10$，従って $q\le25$ となる．

例外辺数は二・四・六のいずれかであるが，三場合は全て排除された．
二例外は，負の大きい reverse を持つ頂点から見た正例外の位置による三場合で，
hub の対向 $D$ 式と正 flow の下界が矛盾する．三高分母の一般分類は不要である．

四例外の再利用できる有限化は，例外の大きい reverse $m=q-1$ がその大きい端の
一つの行だけに現れることを用いる．異符号二項の差を固定した無限方向を
\[
 m_1=h+\max(\delta,0),\qquad m_2=h+\max(-\delta,0),\qquad h\in2\ZZ_{>0}
\]
と書き，$h=2,4$ と全 $h\ge6$ に分ける．最後の半直線では大きい端の
$V=1-R-N$ が必ず非零なので，pair 条件は単位端の node を正例外で零，負例外で二に固定する．
この後の row，node，hub pair，低辺の $D$ 条件は $h$ に依存しない．
従って $h=6$ は無限型全体の厳密な代表であり，分母cutoffではない．
通常二辺が端点非共有の場合の百十二型では候補零，共有の場合の二百四型では候補六となり，
残る無限型は零である．全六候補に
\[
 (q_1,q_2,q_3)=(5,3,3),\quad r=(-4,1,2),\quad x=(-1,2,1)
\]
のcoreが現れ，残る局所条件は $6A^2+17A+14=0$，判別式 $-47$ となる．
これは実数の高高entriesも許さない．

有限化と独立監査には，reverse-first と node-first の二方向がある．
後者は全 $21^4$ node 行列から通常forestのflowを決め，例外分母の行方程式を別に解く．
両方法は六候補の全行列で一致した．次のgraphでは，一方の表現で無限方向が見えにくい場合に
他方から復元し，同じ有限候補へ至るかを調べる価値がある．
全低coreを持つsupportはこれで全分母一に限られるが，一高分母hubへ移すと
pair productsとrow予算が変わるため，補正を導かずに同じ式を適用してはいけない．

全分母一のatlasや一高分母hubとの比較のため，全低 hub 自身の自由度の表示も残す．
$V_{ij}V_{ji}=0$，行和 $-1$ から，各行に負の非零 entry が少なくとも一つある．
非零 support は四，五，六辺であり，出次数列はそれぞれ
$(1,1,1,1)$，$(2,1,1,1)$，および $(3,1,1,1)$ または $(2,2,1,1)$ である．
出次数一の entry は $-1$，出次数二の行は $(t,-1-t)$ なので，
有限個の有向 support と高々二整数の族として扱える．
特定の全 $-1$ 模型に固定せず，正の weights が必要な段階で cycle 条件を課す．
\subsection{部分的な低辺条件だけでは完成しない二つの模型}
\label{ng16:partial-patches}

次の例は，部分的な必要条件の非空性を示す資料である．大域 atlas の存在例ではない．
$K_{2,3}$ の三つの二価 core を cyclic に標識し，canonical 高 ray $A$ が
一方の中心を向く場合を $+$，他方を向く場合を $-$ とする．
符号 $(+,+,-)$ に対し，canonical データを
\[
 (r,s,z,S,T)=(4,4,5,2,2),\quad(5,10,2,3,3),\quad(5,10,2,3,3)
\]
と取る．三低辺の node triples は $(0,3,2),(2,2,1),(0,2,3)$ だが，
同じ順の対向低係数和は $7,-10,5$ であり，相殺しない．
中心二 core も与えていない．本版の $K_{2,3}$ 排除は，中心間の低辺と高辺の clutching を加えて成立する．
具体的な全行列・係数・node の再計算は，保存した v14 の node 検算器に含まれる．

diamond の二価 cores $2,3$ と全低 core $4$ についても，
三つの相互低辺の node 条件と $D$ 相殺を全て満たす部分模型がある．
行列の添字順をそれぞれ $(0,1,3,4),(0,1,2,4),(0,1,2,3)$ として
\[
 P_2=\begin{pmatrix}0&-6&1&-2\\-8&0&-2&3\\6&-3&0&-4\\-3&2&0&0\end{pmatrix},\quad
 P_3=\begin{pmatrix}0&-14&-2&5\\-7&0&2&-4\\-5&4&0&0\\2&-2&-1&0\end{pmatrix},
\]
\[
 P_4=\begin{pmatrix}0&0&-1&0\\2&0&0&-3\\0&-1&0&0\\-3&0&2&0\end{pmatrix},
 \quad q_2=(7,7,1,1),\quad q_3=(11,9,1,1),\quad q_4=(1,1,1,1).
\]
係数は順に
\[
 (-2,0,0,2),\qquad(-1,3,0,-2),\qquad(0,0,-2,2)
\]
であり，低辺 $23,24,34$ の node triples はそれぞれ
$(0,0,5)$，$(4,0,1)$，$(2,3,0)$ である．
全てに正の weights がある．しかし全低 hub の残差
\[
 p^2_{34}+p^3_{24}-p^2_{43}-p^3_{42}=-3
\]
は零でないので，二つの中心 core を加えた大域完成は不可能である．
行和・pair products・正 weights・係数・nodes の厳密検算は
\nolinkurl{diamond_obstruction_verification_v15.json} に保存した．
部分模型の存在を保ちながら，追加する大域条件のどこで矛盾するかを区別するための例である．


\subsection{house 型 \texorpdfstring{$G_{5,5}$}{G5,5} の全排除から再利用するもの}
\label{nh18:house}

$H=\{01,02,03,12,14,34\}$ は，分母上限なしで排除した．
公開算術稿の定理 \texttt{hs18:no-house} が証明を与える．
二価 core $2$ の負の枝は途中で残し，外側 $3,4$ だけを
「高隣接頂点間が低辺」という既知の補題で正の枝へ制限した．
三価 cores $0,1$ の分類は不要であり，$D$ 相殺を使うのは低辺 $23,24$ だけである．
この型の排除と，現在の三型の存在・不存在の問題は区別する．

\paragraph{共通低 ray の $C/D$ 型による分類．}
正の二高分母標準形の低 rays を $C,D$ とすると，高 ray からの
outgoing / reverse entries は
\[
 (C,A):(x,r),\quad(C,B):(-y,-r-1),\quad
 (D,A):(-s,-S),\quad(D,B):(z,T)
\]
である．正の枝では $x\ge2,S\ge2,T\ge5,L=T-S+1\ge3$ であり，
$y,s,z$ はそれぞれの高分母から一を引いた数の真約数である．
共通低 puncture の signed-divisor 接着と reverse 上界四を組み合わせると，
両 $D$ の型は不可能であり，$C/D$ 混在型では
\[
 \begin{array}{c|c}
 \text{通常}&\text{両 shared rays}=A,\quad r_C=S_D\le4\\
 \text{正例外}&\text{shared ray}=A_C,B_D,\quad r_C=z_D=1
 \end{array}
\]
だけが残る．これは別の graph の隣接二価 cores にも再利用できる．
canonical $D$ row の低低 entry は $-L$ であり，零と取り違えない．

\paragraph{最終的に現れる完全な必要族．}
house 型で中央 core $2$ の高 rays を $A=0,B=1$ とすると，node 条件は唯一の族
\[
 \begin{array}{c|c|ccccc}
 \text{core}&(A,B,C,D)&r&s&z&S&T\\\hline
 2&(0,1,4,3)&1&2t-1&t&2&6\\
 3&(0,4,1,2)&1&1&1&2&6
 \end{array}
\]
へ帰着する．$t\ge1$ は任意の整数であり，探索上限ではない．
$D^2_3=t-2,D^3_2=-1$ なので低辺 $23$ の相殺で $t=3$ となる．
core $4$ は $(A,B,C,D)=(3,1,2,0)$，$a=7,r=1$ で
$D_C=-3$ を要求されるが，この固定データの全七模型は
\[
 D_C\in\{4,9,15,18,19,23,40\}
\]
を満たすので不可能である．
固定 $b=7$ にある共鳴族を捨てたのではなく，この場所の共有 ray は
canonical $A$ と先に強制されており，その共鳴族は適用範囲が異なる．

\paragraph{検算と今後の使用範囲．}
\nolinkurl{verify_house_obstruction_v18.py} は，共通低 ray の全十六型，
node の不等式，物理標識付き node 式と無限族の二十五個の多項式恒等式，
固定 $a=7,r=1$ の全約数を確認する．
最後の七模型は，証明で導いた上界を使う別の original-matrix 復元でも一致する．
検算結果は \nolinkurl{house_support_obstruction_verification_v18.json} に保存する．
これは一般の node 不等式による場合分けと完全な有限計算を組み合わせた証明であり，
数値上限以下の分母を試した経験的排除ではない．
house 型自体の探索は終了し，共通低 ray の型分類と低辺係数相殺を残る型へ移す．

% Reusable results only. Detailed proof: gten18:no-support in arithmetic paper.
\subsection{\texorpdfstring{$G_{5,10}$}{G5,10} の排除から再利用する方法}
\label{ngten18:reusable}

$G_{5,10}=\{01,02,03,12,13,24,34\}$ は，分母の上限なしに排除された．
算術稿の定理 \texttt{gten18:no-support} を参照する．
核心は，三高分母 core の正符号が完全 matching となる定理を，
二価 core の二つの低辺の node 予算と組み合わせることである．
対象は指定した五次の算術的接着模型であり，任意の fibration の分類ではない．

頂点 $4$ の二高分母標準形を $(A,B,C,D)=(2,3,0,1)$ とし，
$a=sS+1,b=zT+1,x=sS/r,y=zT/(r+1)$，$\delta=T-S$ とする．
$q_{23}=1$ なので正の枝であり，$\delta\ge2$ である．
低辺 $04,14$ の node 数をそれぞれ $(n_1,n_2,n_3)$，$(m_0,m_2,m_3)$ とすると，
隣接する三高分母 cores の低 rows が
\[
 R^0=(1-n_1,1-r-n_2,r+2-n_3),\quad
 R^1=(\delta+2-m_0,S+1-m_2,1-T-m_3)
\]
に固定される．各 row には正成分がちょうど一つあり，各 node 組の和は五である．
共通低 puncture $(01;4)$ の三つの接着型は次のように処理する．

\begin{enumerate}
\item 通常型では，有限な node 予算と matching 性だけで
$(r,S,T)=(2,3,5)$ が強制される．整数 pair equation は
$(u-2)(v+1)=-3$ に帰着し，唯一の中央模型は
$(r,s,z,S,T)=(2,2,6,3,5)$，$(a,b)=(7,31)$ である．
しかし反対側の core では，$30$ の signed divisor が
$27\pmod{31}$ であることを要求される．その候補 $27,-4$ はどちらも不適格である．
この部分は $D$ 相殺を使わない．

\item 正例外型では $R^0=(1,-1,-1)$ となる．二つの低辺の $D$ 相殺から
$x-y>0$ と $s-z\le1$ が従う．正の枝の不等式は $s-z\ge1$ を与えるので
$s=z+1$ である．反対側の低 column の二つの負成分がともに $-1$ となり，
$q_{12}+q_{13}=q_{01}+2$ を要求するが，奇分母の parity と矛盾する．

\item 負例外型では $q_{01}\in\{3,5\}$，$\delta=2$ となる．
$h=S-r$，$\alpha=sh/r$，$\beta=z(h+1)/(r+1)$ とすると
\[
 \left(\frac{r^2}{h}-\frac{(r+1)^2}{h+1}\right)\alpha\beta
   =\alpha-\beta+1.
\]
第三角の signed-divisor clutching により，$\alpha,\beta>1$ なら括弧が整数で，
その絶対値は一未満となる．零の場合は $h=r^2/(2r+1)$ となり不可能である．
$\alpha=1$ の例外も保持して処理すると，必ず $\beta=1$ となる．
残る equation は $(h+1)z(z-2)=h-1$ であり，中央模型全体が
\[
 r=3,\quad S=4,\quad T=6,\quad z=2,\quad s=3\alpha,\quad
 a=12\alpha+1,\quad b=13
\]
という無限共鳴族へ帰着する．これを探索上限で切らない．
低辺 $14$ の $D$ と奇分母から $\alpha$ は正の偶数となり，
$q_{12}\ge15\alpha+1$ を得る．低辺 $04$ から $q_{02}\ge8\alpha-7$ を得る．
辺 $24$ の node と matching 性は二つの正整数 $M,N$ に
\[
 M+N=24\alpha+1,\quad M>3q_{02},\quad N>q_{12}
\]
を要求するため，$M+N>39\alpha-20\ge24\alpha+10$ と矛盾する．
\end{enumerate}

再利用すべき点は，低辺の有限な node 予算から隣接三高 core の低 row を固定すること，
matching によってその符号を制限すること，そして第三角の signed divisors で
共鳴族を特定してから $D$ と node の不等式で全族を排除することである．
検算の入口は \nolinkurl{verify_g5_10_obstruction_v18.py} であり，
係数ごとの恒等式二十四本と，証明由来の有限な node の場合分けを保持する．
有界分母探索を一般証明の代わりにしていない．

% Reusable v18 notes; no exploration diary.
\subsection{一高分母 hub の有限型復元と \texorpdfstring{$K_4$}{K4} 一葉型の解決}
\label{nl18:leaf-hub}

高分母 support が $K_4$ 上の $0,1,2,3$ に高辺 $04$ を加えた型 $G_{5,12}$ は排除された．
入力は一高分母の完全標準形と，三高分母 core の正符号 graph が完全 matching となる定理である．
全低 hub の $K_4$ 障害を部分 graph へ拡張したものではない．
$G_{5,7},G_{5,9}$ は別の配置解析で排除された．
この $K_4$ 一葉型の証明を，隣接 core の次数が異なる二型へ無条件に転用したわけではない．

\paragraph{再利用する補正公式．}
任意の hub $4$ について $\ell_i=q_{i4}$ とし，
$R_{ij}=p^i_{4j}$，$X_{ij}=p^i_{j4}$，$V_{ij}=p^4_{ij}$，$N_{ij}=n_{i4;j}$ と置けば，
\[
 R_{ij}=\ell_i(1-N_{ij})-V_{ij},\qquad
 R_{ij}X_{ij}=\ell_iq_{ij}-1,\qquad
 \sum_{j\ne i}N_{ij}=5.
\]
対向 $D$ 相殺は
\[
 \sum_{j\ne i}(X_{ij}-R_{ji})
 =\sum_{j\ne i}\ell_jN_{ji}-\sum_j\ell_j-\ell_i
\]
となる．一高 hub では $\ell=(Q,1,1,1)$ である．
高 ray を含む hub pair は $Q-1$ であり，零積へ置き換えない．

\paragraph{葉分母の上界を導く仕組み．}
標準形を $(A,B,C,D)=(0,1,2,3)$ とすると
$u=th$，$v=t(h+1)$，$Q=th(h+1)+1$ である．
特別な $B$ ray の低 row は $R_{10}=2-N_{10}$，$R_{12},R_{13}\le1$ となる．
一正二負条件と clutching から $R_{10}=1,2$ のみが残り，
$q_{01}=2Q-1,Q,2Q+1$ の三族と，$Q=3,5$ における任意奇数 $q_{01}$ の退化族を得る．
$2Q-1$ 族は高 hub 辺の $D$ と他の高辺の node 条件により $Q=3$ に限られる．
$Q>21$ なら $u\ge5,v\ge6$ である．$d=0$ は $R_{31}\ge3$ と
$R_{13}\in\{-1,-2\}$ の反対符号接着に反する．$c=0$ では
core $2$ の唯一の正 entry が $R_{20}$ となり，$|R_{23}|\le4$ である一方，
core $3$ の row sum と一正二負条件は $R_{32}\ge5$ を要求する．
従って $Q\le21$ である．

\paragraph{自由整数と他の分母を切らない復元．}
$Q\le21$ の $(h,t)$ は十五個，四つの node 行はそれぞれ二十一通りである．
$c=0,d=z$ または $d=0,c=z$ とし，$z$ は全整数にわたって残す．
辺 $0i$ では $R=R_{0i},r=R_{i0}$ が $z$ に依存せず，
\[
 (Qr+R-\kappa Rr)q_{0i}=R+r,\qquad -4\le\kappa\le1
\]
から固定分母または全無限族 $q_{0i}=1+\operatorname{lcm}(2,|r|)a$ を得る．
低三角形内では固定 reverse $r$ の相手は $-r$，条件付きの $\operatorname{sgn}(r)$，
または $|r|=1$ における同符号の全正偶数である．
従って $z$ の許容集合は点と偶奇付き半直線の有限和となり，正確に交差できる．

固定 node 配置 $345$，partner 型の組 $1192$ のうち符号条件を満たす $579$ 組は，
全て空交差である．これらは手順内の枝を含む型の数で，charts の同型類数ではない．
\nolinkurl{verify_k4_leaf_obstruction_v18.py} と
\nolinkurl{k4_leaf_obstruction_verification_v18.json} に全 $579$ 個の空交差証明書を保持する．
$Q=3$ における標識重複を残し，$c=d=0$ の二重計数も不存在判定のため許している．
この方法では他の高分母や自由な低低係数に探索上限を置かない．

\subsection{二つの葉付き支持の排除から再利用する方法}
\label{nl19:section}

$G_{5,7}$ と $G_{5,9}$ は，一高完全標準形，三高 matching，
二高の正負二分，共通低 puncture の signed-divisor 分類，node 非負性と対向 $D$ 相殺で排除された．
標準 $K_5$ graph，四直線 cones，marked punctures を保つ五次の整数接着問題に関する結論である．
葉を $4$，その高隣接を $0$ とし，葉の標準 rays を $(0,B,C,D)$ とする．
$Q=th(h+1)+1$，$u=th$，$v=t(h+1)$，$cd=0$ の全整数自由度を保持する．

$G_{5,7}$ では $0$ は三高で，その唯一の低 ray の位置 $m$ が $B,C,D$ の全配置を尽くす．
可能な配置は
\[
 (B,N_{0B}=2,Q\le5),\quad(C,N_{0C}=1,u\le4),\quad
 (C,N_{0C}=2,h=t=1),\quad(D,N_{0D}=1,v\le4)
\]
に限られ，従って $Q\le21$ である．
$G_{5,9}$ では $0$ は四高で，残る三高 core の位置 $J=B,C,D$ を全て残す．
$B$ の特別な reverse は $J=B$ なら $1,2$，他なら $1,2,-2$ だけとなる．
$-2$ の枝は $Q=3$ に限られ，三配置の低低 entry と matching の比較から $Q\le13$ を得る．
いずれの上界も，当該 support の仮定を用いて証明した結果であり，数値探索の上限ではない．

再利用する葉接着の式は
\[
 \ell=(Q,1,1,1),\quad R_{ij}=p^i_{4j},\quad X_{ij}=p^i_{j4},\quad V=P^4,
\]
\[
 R_{ij}=\ell_i(1-N_{ij})-V_{ij},\quad\sum_{j\ne i}N_{ij}=5,\quad
 R_{ij}X_{ij}=\ell_iq_{ij}-1,
\]
\[
 \sum_{j\ne i}X_{ij}+\sum_{j\ne i}V_{ji}=-2\ell_i.
\]
高 ray を含む葉の pair product は $Q-1$ であり，零ではない．
各 node 行は二十一通りだが，それだけでは自由整数も他の分母も有限にならない．
固定 reverse の相手は一点または全偶奇半直線で処理し，中央辺では
\[
 (Qr+R-\kappa Rr)q=R+r
\]
を用いる．係数と右辺がともに零なら，
$q=1+\operatorname{lcm}(2,|r|)a$，$a\ge1$ の全族を残す．

$G_{5,7}$ の残余は $160$ partner 型，中央辺の選択まで含めて $212$ 型である．
中心の $D$ 式に必要な情報をまとめた十三行では，自由分母の七行が $q=1,-1,-3$ を要求し，
残る六行は非零の $D$ 残差を持つ．
$G_{5,9}$ では点 $908$ 型と全半直線 $18$ 型が残る．後者は二高 core の canonical な
第二高分母 $b=3$ を要求して矛盾する．低低 pair の逆 entry は別の低辺の node 予算から
$0,-1,-3,-4$ に限られ，点型から $2596$ 個の $D$ 線形系を得る．
$2563$ 個は有理係数で不整合，残る $33$ 個は $3a+3b=-6$，$a,b\ge1$ を要求する．
各矛盾の元の四式による有理線形結合を証明書に保存した．
自由分母が多数残っても，先に全局所行列を完成させず，$D$ の必要条件だけで排除できることを示す．

\paragraph{小さい正 entry と，局所例の使い分け．}
二高 core の二つの低行の正 entries が同時に二以下になることはない．
canonical には $r,T$ であり，正枝は $T>r+1$，負枝の唯一の小候補は
$(s-4)(z+3)=-6$ と整除性に反する．
この補題は，隣接 matching core から node によって小さい上界が伝わる場合に使える．
一方，三高 core の低行 $(-1,-1,1)$ 自体は不可能ではない．例えば
\[
 q=(1,3,31,7),\qquad
 P=\begin{pmatrix}
 0&-1&-1&1\\-2&0&4&-5\\-30&23&0&-24\\6&-4&-9&0
 \end{pmatrix}
\]
は row/pair equations と正値性を満たす．正二閉路の gains は $6/7$ と $92/93$ である．
葉型の排除には隣接 core と node/$D$ を用いることを忘れない．

完全証明は算術稿の \nolinkurl{g5_7_obstruction_v19.tex} と
\nolinkurl{g5_9_obstruction_v19.tex}，型の全証明書は同名の主検算と結果 JSON にある．
型数は canonical branches の重複を残す手順上の個数であり，同型類数ではない．

\subsection{二つの三価頂点を持つ七高辺型の排除と，再利用する有限化}
\label{ng51120:obstruction}

\[
 G_{5,11}=\{01,02,03,04,12,13,24\}
\]
は，正の整数局所 charts，marked clutching，非負 node，対向 $D$ の相殺の下で存在しない．
二価 cores は $3,4$，三価 cores は $1,2$，四価 core は $0$ である．
この排除には分母および scale の探索上限を用いない．
次の型へ再利用するため，有限化の仮定と無限族の保持を記録する．

\paragraph{向き別の reverse 上界．}
二価 core を $i$，他方の二価 core を $j$，
固有の高近傍を $h$，残る三価頂点を $k$ とする．
canonical order $(A,B,C,D)$ と ray $j$ の行を，列 $(0,h,k)$ の順で書くと
\[
\begin{array}{c|c|c}
\text{型}&(A,B,C,D)&\text{ray }j\text{ の行}\\\hline
\mathrm I&(0,h,j,k)&(r,-r-1,0)\\
\mathrm {II}&(h,0,j,k)&(-r-1,r,0)\\
\mathrm {III}&(0,h,k,j)&(-S,T,-L)\\
\mathrm {IV}&(h,0,k,j)&(T,-S,-L)
\end{array}
\]
となる．共通低 puncture の通常 reverse の絶対値は四以下である．
正例外 $\{1,q-1\}$ と負例外 $\{-1,-(q-1)\}$ を捨てずに比較する．
正の単位 entry $x=sS/r=1$ は，正枝では $S>r$，
負枝では pair equation の符号から一般に排除できる．
これにより I 型は $r\le3$，II 型は $r\le4$，IV 型は $S\le4$ となる．

正の III 型は先に除く必要がある．相手の node と三高 matching により，
両二価 cores が $(r,s,z,S,T)=(1,1,1,2,6)$ に固定される．
外側低辺の $D$ 相殺が $q_{01}+q_{12}=4$ を要求し，両分母三以上と矛盾する．
負の III 型は $w=S=T\le4$ であり，他の負枝も $w\le4$ となる．
$w^2/2<H(H+1)<2w^2$ から $H\le5,r\le9$ を得る．

正枝の相手別の議論は，上界を他の graph へ移す際に確認すべき箇所である．
I--II では相手 $(r',S',T')=(1,2,5)$，
I 側は $(r,S,T)=(3,4,6),(3,5,7)$ に限られる．
II--II では $r=r'=1$，$S,S'\le3$ となり，pair equation が $T,T'<14$ を与える．
II--IV は通常接着なら $T\le5$，負例外なら $L=3$，
$r<S<2r$ から $S\le7,T\le9$ となる．
IV の残る相手は II または負の III であり，それぞれ $T\le6,T\le5$ である．
従って全体は，数学的に証明された有限範囲
\[
 1\le r\le9,\qquad 2\le S\le7,\qquad S\le T\le13
\]
に含まれる．この範囲は全候補を含む上界であって，全ての点が接着可能という意味ではない．

\paragraph{全 scale を保持した点・半直線表示．}
$h=S-r,\delta=T-S$ とすると
\[
 Fsz+As+Bz+C=0,
\]
\[
 F=(2r+1)h-r^2(\delta-1),\quad
 A=h(r+1),\quad B=-r(h+\delta-1),\quad C=r(r+1)
\]
である．$F\ne0$ では
$(Fs+B)(Fz+A)=K=AB-FC$ の全 signed divisors を用いる．
$K=0$ の二直線と $F=0$ の一次方程式は，整数合同式を解いて全半直線として残す．
この graph では book と異なり $L\le4$ を仮定できず，$F=0$ は実際に生じる．
無限族を持つ reverse triples は
\[
\begin{array}{c|c|l}
(r,S,T)&\text{場合}&\text{全正整数 scale}\\\hline
(1,2,6)&F=0&s=2z-1,\ z\ge1\\
(2,6,12)&F=0&s=3a-2,\ z=2a-1,\ a\ge1\\
(2,3,6)&K=0&s=2a,\ z=1,\ a\ge1\\
(3,4,6)&K=0&s=3a,\ z=2,\ a\ge1
\end{array}
\]
である．後二つだけを共鳴として残す手続きでは不完全になる．
整除性と奇数分母条件による剰余類分割の後，一つの二価 core あたり
$90$ 点型と $78$ 全半直線型，計 $168$ 表示を得る．
向き I,II,III,IV の表示数はそれぞれ $57,61,8,42$ である．
これは幾何的模型の個数ではなく，無限族を含む完全なパラメータ表示数である．

\paragraph{最後の低辺と，独立監査の役割．}
低辺 $34$ の node を通る型対は $3592$ 個である．
外側低辺 $14,23$ で各二十一 node 分配を使い，三高の低 row の一正二負，
共通低 puncture の通常・例外接着を課すと $1395$ 配置が残る．
例外接着の要求分母と affine 分母の一致から，
定数不一致 $759$，parameter 非整数 $550$，計 $1309$ 配置を除く．
残る $86$ 配置では低辺 $34$ の残差 $D^3_4+D^4_3$ が
$72$ 配置で常に奇数，$14$ 配置で非零定数となる．
定数の内訳は $8$ が四配置，$28$ が一配置，$-4$ が九配置である．
自由 parameter の上限を設けず，全ての非負整数に対して相殺を排除する．

主検算 \nolinkurl{verify_g5_11_obstruction_v20.py} は
全 affine 行列，node 行，要求分母，残差を証明書に保存する．
独立した \nolinkurl{audit_g5_11_certificates_v20.py} は主生成器を import せず，
canonical row/pair 式と元の node 定義から全 $1395$ 配置を再構成する．
さらに整数性の primitive step を用いる別の生成法で
$90$ 点と $11$ primitive 半直線を得て，
主検算の $78$ 剰余類半直線がその全体を覆うことを有限周期ごとに照合する．
この監査は半直線の最初の有限点だけを照合するものではない．
無上界の結論は解析的 reverse 上界と全整数族の処理の組合せに基づく．

\subsection{外側 core の一つの pair から尺度を復元する}
\label{noc23:section}

低辺の両端だけを先に完成すると，自由尺度が残ることがある．
第三 core の全行列を完成する前でも，その一つの pair equation は必要条件として使える．
これは密な支持の接着へ移植できる今回の主要な道具である．

三つの cores を $i,j_0,k$ とし，$q_{ij_0}=1$，$q_{ik}=a>1$，
$p=p^i_{kj_0}\mid a-1$ とする．低辺の反対側の対応 entry を
$p^{j_0}_{ik}=\varepsilon j$，$p^{j_0}_{ki}=\varepsilon S$，
$\varepsilon=\pm1$，$j,S>0$ と書けば $q_{j_0k}=1+jS$ である．
$u=1-n_{ik;j_0}$，$v=1-n_{j_0k;i}$，$E=au-p$ と置くと，
\[
 \{E(jv-\varepsilon)-aj\}S=a-1-Ev,
 \qquad u,v\in\{-4,-3,-2,-1,0,1\}.
\]
係数が非零なら，各 node 組から正整数尺度は高々一つである．
さらに係数を $A$ とすると
\[
 (1+jS)A=-\varepsilon E-j\ne0
\]
なので，未知分母は固定整数の約数としても求められる．
低辺自身の node 条件 $(a-1)/p+\varepsilon j\in[-4,1]$ を加えると
\[
 j\le a+3,\qquad q_{j_0k}\le6a+2
\]
を得る．この比較は非共鳴時に限る．

共鳴は $E=-\varepsilon j$，$jv=-\varepsilon(a-1)$ で完全に特徴づけられる．
低辺の node 条件まで課した場合，$a\ge7$ で残るのは
\[
 (\varepsilon,p,j,u,v)=(-1,1,a-1,1,1),\quad
 (1,-1,a-1,-1,-1)
\]
だけである．$a=3,5$ には三つの小例外が加わるので，公開稿
\nolinkurl{outer_core_linear_transfer_v23.tex} の全表を使う．
奇数分母，残る pair equations，row sums，$D$ 相殺，正 weights は別に確認する．
単一 pair の共鳴を大域接着の存在と読み替えない．

\nolinkurl{verify_outer_core_transfer_v23.py} は係数恒等式を厳密に確認し，
元の pair 方程式との有限回帰で全共鳴表と低辺を加えた表を照合する．
共鳴の無上界分類は本文の整数整除の証明によるものであり，回帰範囲からの外挿ではない．

% v23 research note: retain reusable equations, scope, and next frontier.
\subsection{通常型の全排除：外側 pair を早く使う}
\label{notes:g51523-ordinary}

\paragraph{確定した新しい結論．}
公開補節 \nolinkurl{g5_15_ordinary_outer_pair_obstruction_v23.tex} は，
標準五次整数接着の $G_{5,15}=K_5-\{23,24\}$ 通常型を全十四配置・全分母で排除する．
通常型の共有分母 $Q$ を増やす探索，非共鳴 $T\ge2$ の通常型 lifting，
通常型の外側四高 core 完成は今後の課題から外す．
第22版の共鳴・$T=1$ 排除は正しい局所補題として保持できるが，通常型の新しい証明には不要である．
第21版の正例外排除と，次節の負例外 $Q=3,5$ の全尺度排除を合わせ，
$G_{5,15}$ 全体が解決した．通常型の専用探索を再開しない．

\paragraph{再利用する三角形の一次式．}
中心と隣接三高 core の間の低辺を $2p$ とする．外側 core $i$ について
\[
 a=q_{2i},\quad c=p^2_{ip},\quad
 p^p_{2i}=-j,\quad p^p_{i2}=-\sigma,\quad q_{ip}=1+j\sigma
\]
と置く．node 変数 $u=1-n_{2i;p}$，$v=1-n_{ip;2}$ はともに $-4,\ldots,1$ であり，
$E=au-c$ とすると
\[
 \{E(jv+1)-aj\}\sigma=a-1-Ev,
 \qquad
 (1+j\sigma)\{jEv+E-aj\}=E-j.
\]
これは未完成の外側 core 内の一つの pair equation そのものである．
係数と定数がともに零となる共鳴は $E=j$，$v=(a-1)/j$ に限られる．
従って $a-1>j$ なら共鳴はなく，$q_{ip}\mid E-j\ne0$ から一様有限化できる．
中心の全体を完成してから外側 core を調べる順序は必須ではない．
他の支持でも，既知の低辺の両端から外側 core の必要 pair を先に課す方法を試す価値がある．
ただし $a-1\le j$ の共鳴は一般に除いてはならない．

\paragraph{通常型での有限表と最後の矛盾．}
外側二 core の尺度集合は，正側の順序を保って次の通りである：
\[
\begin{array}{c|c|c}
 (\text{中心},R^p,d)&S&T\\\hline
 (M_1,(-1,-4,4),5)&2&5\\
 (M_1,(-1,-1,1),-9)&2&\varnothing\\
 (M_2,(-1,-2,2),2)&2&1\\
 (M_2,(-2,-1,2),2)&\varnothing&4\\
 (M_3,(-1,-2,2),-7)&\varnothing&1\\
 (M_4,(-1,-2,2),9)&\varnothing&13
\end{array}
\]
第三行は $U=S+T-d-1=0$ で不可能．第一行だけが $U=1$ を残すが，
高分母 $(3,21,5)$ では元の pair・行和から
$HK=62$，$(H+1)M=14$，$M+N=6$ を得る．
最初の二つの整除性は $H=1,M=7$ を強制し，$N>0$ と矛盾する．
この表の全探索範囲は各外側入力の三十六 node 組から導かれ，任意の分母 cutoff ではない．
二つの負成分を交換した $M_2$ の三配置も独立に含めている．

\paragraph{局所分類との区別．}
通常型が消えたことから，五つの固定 $(R,D_0)$ の孤立局所問題がすべて不存在になったとはいえない．
例えば $R=(-1,-2,2)$，$D_0=2$，$(S,T,U)=(6,3,6)$ の正整数局所模型は実在し，
中心と外側 pair の組合せに接着できない．
独立した局所分類は他の支持に役立つ場合に限り続ける．
Schur の $k\ge0$ と $R=(-1,-m,m)$ の順を同時に固定しないという第22版の注意も引き継ぐ．
今回の証明は Schur の向きを使用していない．

\paragraph{検算の入口．}
\nolinkurl{verify_g5_15_ordinary_outer_pair_v23.py} は旧版 import なしで
四中心，十四配置，六 profile，十一個の相異なる外側入力を復元する．
node 直接列挙と約数列挙を独立に照合し，元の外側 pair equation を各残候補で確認する．
最後の一点は正約数の全検査と，二次残余 $2H^2-35H+124$ の判別式 $233$ の非平方性の
両方で確認する．後者は独立な補助確認であり，主証明は短い整除矛盾で完結する．
証明・ノート・検証器の hash と検証範囲を結果 JSON に記録する．

\subsection{負例外は全排除済み：再利用する外側 pair と resultant}
\label{notes:negative-exception-v23}

第23版で $G_{5,15}$ の負例外 $Q=3,5$ の二 template は全尺度で排除された．
通常型・正例外の排除と合わせて $G_{5,15}$ 全体が解決し，現行支持は三型・二十六標識となる．
この型の大きい尺度や小共有分母を再探索する必要はない．
指定した標準 $K_5$ 接着の障害であり，任意の torus fibration の不存在には拡張しない．

\paragraph{未完成の第三 core を先に使う．}
低辺の両端を完成してから外側 core を調べる順序は，無限尺度を余計に残す場合がある．
外側 core の一つの pair equation と二つの node だけを先に課せる．
中心の分母 $a$，entry $p$，隣接低成分 $\varepsilon j,\varepsilon A$ に対して
\[
 E=au-p,\quad -4\le u,v\le1,\qquad
 A\{E(jv-\varepsilon)-aj\}=a-1-Ev
\]
である．共鳴は係数と定数の両方が零になる条件から完全に分類する．
分母 $q=1+jA$ に対する $q\{E(jv-\varepsilon)-aj\}=-\varepsilon E-j$ も再利用できる．

\paragraph{今回の定数シフト．}
中心を $(r,s,z,S,T,x,y)$ とし，$J$ の低 column を $(-A_J,B_J,-1)$，
$K$ のそれを $(-A_K,B_K,-(Q-1))$ とする．全有効候補は
\[
\begin{array}{c|c|ccc|ccc}
&Q&j&A_J&B_J&k&A_K&B_K\\\hline
N_\pm&3&r+\varepsilon&x+\varepsilon&y+\varepsilon&w+\eta&z+\eta&s+\eta\\
P_\pm&5&r-1&x-1&y+1&S+1&z-1&s+1
\end{array}
\]
へ落ちた．$\varepsilon,\eta=\pm1$，$P$ は正 $L=3$ の枝である．
奇数分母は全族で $\ell$ 奇数・$t$ 偶数を強制する．
これは元の局所族の parity が自動であることとは別であり，接着後の新しい三高分母へ課した条件である．
十三の共鳴 $\ell=0$ は自由な $t$ を全て保持した後，隣接高分母六として排除される．
七の正 $L=4$ 共鳴は，この template の node 条件 $L+2\le5$ に入らない．

\paragraph{再利用できる単位 reverse column の残余式．}
一般に低 row が $(-j,j+m-1,-m)$，低 column が $(-A,B,-1)$，共有分母が $Q=m+1$ なら，
共有高 row は $(X,-X-m,-1)$ である．元の残余は $v=j-X$ を用いて
\[
 [X(X+m)-Qv(v-1)]AB+XvA+(X+m)(1-v)B+X(X+m)=0
\]
となる．$D_{\rm low}=A-B$ なので，既知の隣接 $D$ を直接入れられる．
この式自体は $Q=3,5$ に限らない．ただし必要な整除性・正値性を捨てて十分条件にしてはならない．

\paragraph{二つの二次式を同時に消す方法．}
共有高 node の整数 $n=0,1,2,3$ は $X-Y=Q+1-Qn$ を与える．
二つの元残余を $f=aX^2+bX+c$，$g=dX^2+eX+f_0$ とすると，
共通根には $(af_0-cd)^2-(ae-bd)(bf_0-ce)=0$ が必要である．
これは先頭係数で割らない整数恒等式なので，一次退化も保たれる．
今回の四十場合は，三つの全無限領域と四つの端点に分けた係数符号証明で排除した．
唯一の零 resultant は $N_-,\varepsilon=\eta=-1,n=2,\ell=1,t=2$ で，$B_J=0$ となる．
resultant の零を局所模型の存在と取り違えない．

\paragraph{復元資料．}
公開証明は \nolinkurl{g5_15_negative_exception_obstruction_v23.tex}，
全係数証明は \nolinkurl{verify_g5_15_negative_exception_v23.py} と対応する JSON にある．
検算は標準ライブラリのみで，外側 node の千三百六十八個の記号的場合分け，
四十 resultant，百二十の全係数符号，四端点での百六十値を確認する．
元の正整数族の任意 cutoff を設けた検算ではない．
残る支持へこの手順を移す際は，低辺・中心の種類・外側 node 範囲を改めて導く．

\subsection{Book は排除済み：再利用する整数剛性と尺度公式}
\label{nb22:section}

$G_{5,13}=\{01,02,03,04,12,13,14\}$ の二価 cores の低三角形は存在しない．
第21版で残った二つの一負枝 sector を，第22版では探索上限を増やさず排除した．
対象は引き続き標準 $K_5$ 配置，正の整数局所 charts，奇数高分母，
非負 node・総和五，対向低 $D$ 相殺である．
任意の Lagrangian torus fibration に対する一般的な不存在定理にはしていない．
今後 book の大きい分母を探索する必要はない．

\paragraph{証明の依存関係と再開時の確認点．}
第21版の定理 \texttt{bs21:two-sectors} は，既存の全尺度有限復元
$2\le w\le134$ を含めて，仮想的な book を次へ帰着している：
\[
 \begin{array}{c|c|c}
 &\text{node 周期}&(L_{P_1},L_{P_2})\\\hline
 \mathrm I&N^+\to P_1^+\to P_2^+\to N^+&(3,3)\\
 \mathrm {II}&N^+\to P_1^+\to P_2^-\to N^+&(4,3)
 \end{array}
\]
どちらも $\tau=+$，$S_N=T_N=w\ge135$，$H=r_N-w\ge w$ を満たす．
新しい証明は \nolinkurl{paper/book_complete_obstruction_v22.tex}，
補助的な厳密係数検査は
\nolinkurl{verification/verify_book_complete_obstruction_v22.py} と対応する JSON にある．
この新検査器は旧 generator を import せず，百五十二個の多項式恒等式，
型 I の全十六 offset，型 II の parity 障害，全六族の尺度公式と
高分母の自動的な奇数性を検査する．
既存の第21版の有限証明書に代わるものではない．

\paragraph{小さい $L$ の正枝：厳密な reverse data．}
一般正枝で
\[
 \begin{gathered}
 \delta=T-S\in\{2,3\},\quad d=\delta-1,\quad h=S-r,\\
 \alpha=sh/r,\quad\beta=z(h+d)/(r+1),\quad F=(2r+1)h-r^2d.
 \end{gathered}
\]
と置く．$\alpha,\beta$ は正整数であり，$F>0$ なら
\[
 0<\frac{F\alpha}{h(h+d)}
 =\frac{r(2\alpha-ds)+\alpha}{h+d}<1.
\]
$d=1$ では非負整数 $2\alpha-s$，$d=2$ では非負整数 $\alpha-s$ が
一以上であると右辺が一以上になる．従ってそれらは零であり，
\[
 \begin{array}{c|ccc}
 L&r&S&T\\\hline
 3&2R&3R&3R+2\\
 4&r&2r&2r+3
 \end{array}
\]
に限られる．これは book の接合に依存しない局所剛性なので，
他の支持の局所模型にも使える．

\paragraph{同じ正枝の全尺度：将来の接合への公式．}
$L=3$ の場合，$R\ge1$ を固定する．
\[
 \ell\mid R+2,\quad1\le\ell\le R,\quad
 t=\frac{R+2}{\ell}-1,\qquad
 s=2(R+1-\ell),\quad z=(2R+1)t
\]
が全尺度を与える．$R=\ell(t+1)-2$ なので
$Rt=\ell t(t+1)-2t$ は偶数であり，高分母の奇数性も自動である．
$L=4$ の場合，$r\ge1$ を固定する．
\[
 \ell\mid r+3,\quad1\le\ell\le r+1,\quad
 t=\frac{r+3}{\ell}-1,\qquad
 s=r+2-\ell,\quad z=(r+1)t
\]
が全尺度を与える．$(r+1)t=\ell t(t+1)-2t$ は偶数なので，
こちらも高分母の奇数性は自動である．
これらは必要条件だけではない．表示式を代入すると整数行列の全 row 和と
全六 pair 方程式を満たし，正枝の strict chart 条件も既知の正の weights によって満たされる．

必要性は，$L=3$ で $\alpha=s/2$，$E_+=1/(R+1)$，
$\beta=z(R+1)/(2R+1)$ を使い，$\gcd(R+1,2R+1)=1$ から
$\beta=(R+1)t$ と書けば
\[
 (t+1)(R+1-\alpha)=R+2
\]
となることによる．$L=4$ も $\alpha=s$，$E_+=1/(r+2)$ と
$\gcd(r+2,r+1)=1$ から
$(t+1)(r+2-s)=r+3$ を得る．
従って共鳴自由尺度の取りこぼしはない．

$p=-D_D>0$，$q=D_C>0$ と書けば，この表示は低 $D$ の大きさも与える：
\[
 \begin{array}{c|c|c|c}
 L&p&q&\text{線形関係}\\\hline
 3&2\ell t^2-3t+3&3\ell t^2-4t+5&2q-3p=t+1\\
 4&\ell t^2-2t+2&2\ell t^2-3t+5&q-2p=t+1
 \end{array}
\]
この公式は別の低 cycle を扱う際に，符号だけでなく大きさを接合する入口になる．

\paragraph{正枝の $F<0$ 側も剛性を持つ．}
$\rho=r+1$，$k=h+d$ と書くと
\[
 F=(2\rho-1)k-\rho^2d,\qquad
 0<E_-\beta=\frac{\rho(dz-2\beta)+\beta}{h}
 =1-\frac{\beta-1}{\alpha}\le1.
\]
$F<0$ と整数性から，$d=1$ では $k\le\rho/2$，
$d=2$ では $k\le\rho$ である．従って非負整数
$z-2\beta$ または $z-\beta$ が一以上なら，表示式の分子が分母を超える．
それらは零なので，
\[
 \begin{array}{c|ccc}
 L&r&S&T\\\hline
 3&2R-1&3R-2&3R\quad(R\ge2)\\
 4&r&2r-1&2r+2\quad(r\ge2)
 \end{array}
\]
に限られる．$F=0$ は $\gcd(r^2,2r+1)=1$ から
$2r+1\mid d\in\{1,2\}$ を要求するので不可能である．

この二族の全尺度は $\ell\ge0$，$t\ge1$ によって
\[
 \begin{array}{c|c|c|c}
 L&\text{reverse parameter}&s&z\\\hline
 3&R=\ell(t+1)+2&(2R-1)t&2(\ell t+1)\\
 4&r=\ell(t+1)+2&rt&\ell t+1
 \end{array}
\]
となる．必要性は，$L=3$ で $\alpha=(R-1)t$，$\beta=z/2$ から
$(t+1)(R-1-\beta)=R-2$，$L=4$ で $\alpha=(r-1)t$，$\beta=z$ から
$(t+1)(r-1-z)=r-2$ を得ることによる．逆向きも行列へ直接代入できる．
どちらも $Rt$ または $rt$ が $\ell t(t+1)+2t$ で偶数となるので，高分母は自動的に奇数である．

$\ell=0$ は任意の $t\ge1$ を持つ既知の共鳴族
$(r,s,z,S,T)=(3,3t,2,4,6),(2,2t,1,3,6)$ そのものである．
それ以外の固定 reverse data の尺度は $R-2$ または $r-2$ の正約数によって有限に復元できる．
従って「固定 reverse data の尺度は全て有限」と共鳴族まで含めて述べてはいけない．

低係数は
\[
 \begin{array}{c|cc|c}
 L&D_C&D_D&\text{関係}\\\hline
 3&5-3\ell t^2-4t&2\ell t^2+3t-3&2D_C+3D_D=t+1\\
 4&5-2\ell t^2-3t&\ell t^2+2t-2&D_C+2D_D=t+1
 \end{array}
\]
である．$L=3$ の $(\ell,t)=(0,1)$ は $(D_C,D_D)=(1,0)$，
$L=4$ の $(0,1)$ は $(2,0)$，$(1,1)$ は $(0,1)$ を与える．
最後の模型は $(r,s,z,S,T)=(4,4,2,7,10)$ である．
この三点を除くと $D_C<0<D_D$ である．
以上により，負枝の二族と合わせた全六族が $L\in\{1,3,4\}$ の二高模型を尽くす．
公開証明は \nolinkurl{paper/two_high_rigid_scale_families_v22.tex} にまとめてある．

\paragraph{負枝の完全な reverse 分類．}
負枝で
\[
 E=1-\frac{w^2}{H(H+1)},\quad
 c=\frac{sH}{w+H}\ge2,\quad
 e=\frac{z(H+1)}{w+H+1}\ge1
\]
と置くと $Ece=e-c+1$ である．$H^2<H(H+1)<(H+1)^2$ より $E\ne0$ である．
$E>0$ で $H>w$ と仮定すると，
正整数 $u=s-c=cw/H$ は $u\le c-1$ なので
\[
 1>Ec=c-\frac{u^2}{c}\frac{H}{H+1}
 >c-\frac{u^2}{c}\ge2-\frac1c>1
\]
と矛盾する．従って $E>0$ では必ず $H=w$，$r=2w$ である．
$E<0$ で $H\le w-2$ と仮定すると，
正整数 $v=z-e=ew/(H+1)$ は $v\ge e+1$ なので
\[
 1>(-E)e=\frac{v^2}{e}\frac{H+1}{H}-e
 >\frac{v^2}{e}-e\ge2+\frac1e>1
\]
と矛盾する．従って $E<0$ では必ず $H=w-1$，$r=2w-1$ である．
この分類は低辺の接合や book に依存しない．

\paragraph{負枝の全尺度と低 $D$．}
どちらの符号でも $S=T=w$，$\ell\mid w$，$t=w/\ell-1\ge1$ と書ける．
全尺度は次の二族である：
\[
 \begin{array}{c|c|c|c|c}
  &r&\ell\text{ の範囲}&s&z\\\hline
  E>0&2w&1\le\ell\le w-1&2(w+1-\ell)&(2w+1)t\\
  E<0&2w-1&1\le\ell\le w-2&(2w-1)t&2(w-1-\ell)
 \end{array}
\]
両族とも $wt=\ell t(t+1)$ は偶数なので，高分母の奇数性は自動である．
必要性は，$E>0$ で $\gcd(w+1,2w+1)=1$ から
$e=(w+1)t$，$\ell=w+1-c$ とおけば $(t+1)\ell=w$ となり，
$E<0$ で $\gcd(w-1,2w-1)=1$ から
$c=(w-1)t$，$\ell=w-1-e$ とおけば同じ因数式になることによる．
十分性は行列へ直接代入して全 row 和と pair 方程式を確認できる．
低 $D$ は
\[
 \begin{array}{c|cc}
  &D_C&D_D\\\hline
  E>0&\ell t^2-1&-2\ell t^2-t+1\\
  E<0&-\ell t^2-1&2\ell t^2-t+1
 \end{array}
\]
であり，両方で $-D_D-2D_C=t+1$ となる．
零低 $D$ は $E>0$，$\ell=t=1$，$w=2$ の唯一の模型
$(r,s,z,S,T)=(4,4,5,2,2)$ だけである．
以前の一般 factorization は正しいが，負枝だけを生成する場合には，
$w,H$ の二重ループをこの二つの約数族へ置き換えられる．
既存証明書を修正しなくても，今後の局所・接合計算でこの短い完全表示を使える．

\paragraph{二つの sector を閉じる式．}
型 I では $T_i=(3r_i+4)/2$ と node offset から
\[
 4H=5w+\kappa,\quad
 \kappa=20-4t_1-6t_2-9t_N\in[-18,19],
 \quad t_1,t_2\in\{1,2\},\quad t_N\in\{-1,0,1,2\}.
\]
この sector では $E>0$，従って新しい負枝剛性から $H=w$ である．
代入すると $w=-\kappa\le18$ となり，第21版の $w\ge135$ に反する．

型 II では $P_1$ の剛性と最初の node が $r_N=2r_1+1$ を与える一方，
負枝の剛性は $r_N=2w$ を与える．偶奇の矛盾であり，この最後の排除には
第21版の下限 $w\ge135$ さえ不要である．

高中心二 core の四高分母条件や，高中心間の denominator order を使わずに
book を閉じられたため，その専用完成問題は今後の研究課題から外す．
今後は残った支持に対する障害・構成を優先し，上記の局所尺度公式を必要に応じて移植する．

\subsection{残る三型と再開する具体的な入口}
\label{ns23:support-agenda}\label{ns22:support-agenda}\label{ns20:support-agenda}

再開時は最初に \nolinkurl{00_CONTINUATION_POLICY_ja.txt} を読む．
現在の高分母支持の必要候補は\emph{三型・二十六標識}である．
全候補は五頂点を使い，三点連結，最小次数三以上，八高辺以上，$b_1\ge4$ となる．
単位枝は高々二辺の matching であり，同じ core に二本が接すれば全分母一となる．
補 graph の孤立頂点を省略すると，候補は次の通りである．
\begin{center}\small
\begin{tabular}{c|c|c|c|r}
型&高辺数&次数列&低辺補 graph&標識数\\\hline
$G_{5,14}$&8&$(43333)$&$2K_2$&15\\
$G_{5,16}$&9&$(44433)$&$K_2$&10\\
$G_{5,17}$&10&$(44444)$&$\varnothing$&1
\end{tabular}
\end{center}
識別子は詰め直さない．検算は \nolinkurl{verify_support_agenda_v23.py} である．
候補数は整数接着，atlas，幾何的非同型類の存在数ではない．

\paragraph{高分母全体の原始性．}
非空 $H$ には八辺以上ある．全高分母を割る素数があるとすると，素数支持が $H$ と一致するため，
密な素数支持の定理はその素数を五，支持を $K_{2,2,1}$ に強制する．
補 graph の二辺は実際には単位分母であるが，合同分類はその和を零 modulo five にするため矛盾する．
従って高分母全体の最大公約数は一である．

二本の単位枝がある場合には任意七本の分母の最大公約数も一となり，
$G_{5,14}$ で五の倍数の分母は高々六本である．
密な $G_5=K_{2,2,1}$ と一本の実単位枝が共存する場合は $H=G_{5,16}$ に限られ，
五進支持の外のもう一本は単位枝と頂点を共有せず，分母 $q\equiv9\pmod{10}$，$q\ge9$ となる．
合同条件だけで整数 lifting を保証してはいない．

\paragraph{優先候補：二本の互いに素な低辺．}
Book と $G_{5,15}$ は全排除済みである．専用の大尺度探索や外側四高 core 完成は再開しない．
$G_{5,14}$ の二本の低辺には，三高 core 同士が接する．
第三 core の単一 pair による線形復元を両低辺に適用し，
非共鳴の分母比較と共鳴族を分けて全体の閉路条件を比較する．
固定中心四点への還元は $G_{5,15}$ の特別な構造だったため，そのまま移してはならない．
二高局所族，外側尺度表，二つの二次残余の resultant という手順のうち，
どの部分が新しい支持でも使えるかを明示する．

\paragraph{並行候補：三高同時非負性の残る格子領域．}
固定 $(s,t)$ ごとの全係数復元に，新しい四領域の排除を加える．
一般の $|\delta|\ge2$ では，尺度の正値性が許す $x,n$ の幅を先に制限する．
$\delta=0,s<t$ を標識交換だけで除けると仮定しない．canonical な $k$ の向きも変わるためである．
$\delta=1,x=0$ は $s=y=1,z\ge2$ の軸となり，固定 $(b,z)$ ごとの非零二次式で完全復元できる．
退化面 $a=0$ の大部分は固定 $s$ で有限化でき，残る $x=0,R_0=1$ の枝は
$s^2\mid2(b-1)$ を要求する．この二つの具体的な残余は別々に進められる．
反例が現れれば，全ての原方程式・奇数性・六符号・正 weights を確認して構成へ転換する．

\paragraph{より密な支持と幾何への転換．}
$G_{5,16},G_{5,17}$ では四高 core の符号 graph と元の pair・行和条件が重要になる．
三高 matching 定理を四高へ適用しない．$K_4$ を部分 graph として含むことだけでも排除できない．
整数候補が残れば全 core の weights，basin areas，integer framings，作用切片へ進む．
五次・六次 atlas の同値，不変量，別の実際の退化，open-corridor surgery への転換も，
研究全体の面白さとインパクトに応じて判断する．

